\documentclass[11pt,a4paper]{article}
\usepackage[T1]{fontenc}
\usepackage{iftex}
\ifPDFTeX\usepackage[utf8]{inputenc}\fi
\usepackage{lmodern}
\usepackage[margin=25mm,headheight=15pt]{geometry}
\usepackage{amsmath,amssymb,amsthm,mathtools}
\usepackage{booktabs,array,tabularx,enumitem,graphicx,float,needspace}
\usepackage{microtype,xcolor,fancyhdr}
\usepackage[hyphens]{url}
\usepackage[colorlinks=true,linkcolor=blue!45!black,citecolor=blue!45!black,urlcolor=blue!55!black]{hyperref}
\hypersetup{pdftitle={A general entropy-tail law for adaptive rare-event estimation},pdfauthor={Consolidated working manuscript},pdfsubject={Standalone proof chain, entropy-tail selection, and exact finite certificates}}
\setlength{\parindent}{0pt}
\setlength{\parskip}{0.4em}
\setlength{\emergencystretch}{2em}
\setlist{itemsep=0.2em,topsep=0.3em,leftmargin=1.5em}
\pagestyle{fancy}
\fancyhf{}
\fancyhead[L]{\small An entropy--tail law for rare-event estimation}
\fancyhead[R]{\small Referee-audit revision, 29 October 2026}
\fancyfoot[C]{\thepage}
\renewcommand{\headrulewidth}{0.3pt}
\newtheorem{theorem}{Theorem}[section]
\newtheorem{lemma}[theorem]{Lemma}
\newtheorem{proposition}[theorem]{Proposition}
\newtheorem{corollary}[theorem]{Corollary}
\theoremstyle{definition}
\newtheorem{definition}[theorem]{Definition}
\theoremstyle{remark}
\newtheorem{remark}[theorem]{Remark}
\newcommand{\E}{\mathbb E}
\newcommand{\Pp}{\mathbb P}
\newcommand{\R}{\mathbb R}
\newcommand{\1}{\mathbf 1}
\DeclareMathOperator{\Var}{Var}
\DeclareMathOperator{\KL}{KL}
\DeclareMathOperator{\TV}{TV}
\DeclareMathOperator{\Ent}{Ent}
\DeclareMathOperator{\Aff}{Aff}
\DeclareMathOperator{\kl}{kl}
\title{\textbf{A general entropy--tail law for\protect\\adaptive rare-event estimation}\protect\\[0.4em]
\large Distributional characterization and certified stopping programs}
\author{}
\date{29 October 2026\protect\\Revised after consolidated referee audit}
\begin{document}
\maketitle
\thispagestyle{empty}
\begin{abstract}
A population unit has a hidden positive magnitude drawn from a known law
and a hidden sign whose conditional probability may vary arbitrarily with
that magnitude, subject to a global rarity constraint. Repeated observations
are Bernoulli, and every observation is charged. We characterize the minimax
root-mean-square error under a uniform expected budget, allowing arbitrary
creation, revisiting, interleaving, and stopping. Up to universal constants,
the risk is an entropy--tail selection functional: a selected part of the
law pays its rare-event inverse-square scale entropy, and the remainder
pays a truncated inverse-gap cost. The result holds jointly for every
magnitude law, every allowed rarity, and every positive budget without
moment or regularity assumptions. The upper bound uses an observable
classifier race; the lower bound combines optimal selection, entropy
tiers, scale-sensitive annular saturation, and additive localization.
We also derive an exact attained four-prior dual of the robust stopping-flow
program, a factor-sixteen comparison with minimax risk, and convergent
latent-reveal brackets that certify a sufficient finite cap. An odd-moment
argument identifies the fixed-cap bias floor. All specialized proof inputs
are included, with their methodological antecedents credited. Exact rational
certificates accompany eight finite numerical instances.
\end{abstract}
\medskip
\fbox{\begin{minipage}{0.94\textwidth}
\textbf{Review status.} This revision incorporates the consolidated referee audit, including
explicit genie-convergence scope, strengthened ordinary-mass and sign-entropy
bounds, and clarifications of the adaptive proof interfaces. Its finite
numerical bounds have exact rational certificates.
External mathematical refereeing and proof-assistant verification of the
general theorems remain outstanding. The comparison with prior work is
explicit and targeted, without a publication-priority certification.
\end{minipage}}
\newpage
\begingroup\small\setlength{\parskip}{0pt}\tableofcontents\endgroup
\newpage
\section{Model and main results}\label{sec:overview}

\subsection{The observable experiment}

Let $\mu$ be a known probability measure on $(0,1/2]$.
Each fresh unit independently samples a hidden magnitude $D\sim\mu$
and a hidden sign $S\mid D=d\sim\operatorname{Ber}(q(d))$.
Conditional on $(D,S)$, its observations are iid Bernoulli with mean
$1/2+(2S-1)D$. Different units have independent latent states and
observation tapes. The unknown measurable allocation belongs to
\[
 \mathcal Q_\rho=
 \{q\in L^\infty(\mu):0\le q\le1,\ \theta(q):=\textstyle\int q\,d\mu\le\rho\},
 \qquad 0<\rho\le1/4.
\]
The target is the positive-sign proportion $\theta(q)$.

An algorithm knows $\mu,\rho,B$, but observes neither $D$ nor $S$.
Using an independent random seed and the observed history, it may request
an observation from a fresh or previously observed unit, interleave
units, stop, and return a measurable real estimate. Every requested
observation costs one, including the first on a unit. Units created but
never observed have no statistical effect and may be sampled only when
their first observation is requested. If $T_A$ is total work, define
\begin{equation}\label{eq:model-risk}
 R_{\mu,\rho}(B)=
 \inf_{\substack{A:\ \sup_{q\in\mathcal Q_\rho}\E_qT_A\le B}}
 \ \sup_{q\in\mathcal Q_\rho}
       \{\E_q(\widehat\theta_A-\theta(q))^2\}^{1/2},
 \qquad B>0.
\end{equation}
This is a uniform expected-budget criterion. Stopping may be unbounded;
finite expected work gives almost-sure finite termination under every
admissible population. Estimates may be clipped to $[0,\rho]$ without
increasing risk. In particular, the constant $\rho/2$ has risk at most
$\rho/2$, and randomization permits meaningful procedures even for $B<1$.

\subsection{An entropy--tail selection law}

For a submeasure $\sigma\le\mu$, let $\Gamma_\sigma$ be its homogeneous
dyadic inverse-square cost entropy. With bands
$I_j=(2^{-j-2},2^{-j-1}]$, put
$a_j=\int_{I_j}d^{-2}\,\sigma(dd)$ and $A=\sum_j a_j$.
When $0<A<\infty$,
\[
 \Gamma_\sigma=A[1+\Ent((a_j/A)_j)].
\]
The finite-fragment definition in Section~\ref{sec:additive} also covers
infinite costs and entropies; $\Gamma_0=0$. All logarithms are natural.
Our main result is the following uniform comparison:
\begin{equation}\label{eq:overview-main}
 \boxed{\qquad R_{\mu,\rho}(B)\asymp\mathcal K_{\mu,\rho}(B),\qquad}
\end{equation}
where
\begin{equation}\label{eq:overview-selection}
 \mathcal K_{\mu,\rho}(B)=
 \min\left\{\rho,\inf_{E\ {\rm measurable}}
 \left[\sqrt{\frac{\rho\Gamma_{\mu|E}}{B}}
 +\int_{E^c}\min\left\{1,\frac1{d\sqrt B}\right\}\mu(dd)\right]\right\}.
\end{equation}
Theorem~\ref{thm:entropy-tail-law} states the result formally.
Both comparison constants are universal, simultaneously for every
probability law on $(0,1/2]$, every $0<\rho\le1/4$, and every $B>0$.
No density, regular variation, positive lower support bound, finite
inverse moment, or finite full-law entropy is assumed.

The selected set pays rare-event cost entropy. Its complement pays an
ordinary inverse-gap cost for $d>B^{-1/2}$ and only its probability mass
below that scale. Equivalently, one minimizes
\[
 \zeta((0,1/2])+
 \frac{\sqrt{\rho\Gamma_\sigma}+\int d^{-1}\tau(dd)}{\sqrt B}
 \quad\text{over }\sigma+\tau+\zeta=\mu,
\]
and caps by $\rho$. Hard measurable assignments suffice; the optimizer
need not be an interval. The split is analytical: the estimator never
observes a unit's membership in a selected component.

\subsection{Proof structure}

The upper construction races a likelihood bank against an ordinary
capped classifier on independent observation streams of the same latent
unit. Its null likelihood uses the full law $\mu$. A mixed-work bound
and finite Poisson multilevel allocation give the three-piece upper
bound in Section~\ref{sec:additive}.

The lower proof in Section~\ref{sec:matching-lower} begins with an optimal
finite-band selection. Its ordinary and unresolved pieces are controlled
by a variational null-work modulus. Optimality constrains the gaps of
the retained bands. Grouping comparable costs and trimming a controlled
number of the smallest gaps produces annular components with admissible
conditional rarities and active mass-saturation constraints. Allocating
rarity proportionally to retained entropy makes their lower bounds add
at every budget. Finite-prefix approximation covers laws accumulating
at zero.

All specialized ingredients used in these arguments are proved in this
manuscript. The appendices include the ordinary classifier, weighted
likelihood bank, homogeneous work sum, multilevel allocation, weighted
sign entropy, scale-sensitive gate, null-work transfer, and additive
localization. Standard concentration, minimax, and convex-duality tools
are cited where used. Section~\ref{sec:related-work} distinguishes these
ingredients and their antecedents from the contribution of the joint
selection law.

\subsection{Quantitative stopping certificates}

A complementary representation retains more quantitative information.
The robust cap-$N$ stopping-flow program $P_N(B)$ has exact attained dual
\[
 P_N(B)^2=
 \max_{q_L,q_H,q_v,q_w\in\mathcal Q_\rho}
 \frac{[\theta(q_H)-\theta(q_L)]_+^2}
 {4+B K_N(q_L,q_H;q_v,q_w)}.
\]
Here $K_N$ is optimal stopped signed-likelihood information relative to
$q_v$ per expected observation under $q_w$.
The four priors have separate endpoint, variance, and work roles.
Theorem~\ref{thm:exact-dual} gives the identity and
Proposition~\ref{prop:dual-bellman} gives its Bellman recursion.

Reference absorption and affinity testing give the explicit comparison
\begin{equation}\label{eq:overview-program}
 R_{\mu,\rho}(B)\le P_{N_*}(B)\le16R_{\mu,\rho}(B),
 \qquad
 N_*=\left\lceil64\max\{B,\rho B^2\}\right\rceil.
\end{equation}
The factor sixteen is a program guarantee, distinct from the conservative
constants in \eqref{eq:overview-main}.

A strengthened lower experiment reveals a unit's latent state at its cap.
Its dual value $G_N$ obeys
\[
 0<G_N\uparrow P_\infty,\qquad
 P_N\downarrow P_\infty,\qquad
 G_N/8\le R_{\mu,\rho}(B)\le P_N.
\]
Thus a verified feasible upper policy $U$ and lower dual value $L$ with
$U\le2L$ certify a factor-sixteen approximation. For every fixed input,
sufficiently accurate program solutions eventually meet that criterion.
No uniform convergence rate or bit-complexity claim is made.
Finite-support inputs give genuine finite conic programs.
The eight rational examples in Appendix~\ref{sec:computation} pass at
caps $1$, $4$, or $8$; independent exact checkers verify their feasibility.

The fixed-cap program also has an exact identification threshold:
for $m$ distinct positive magnitudes its infinite-budget bias floor
vanishes exactly when $N\ge2m-1$. Infinite-support laws have a positive
floor at every fixed finite cap
(Theorem~\ref{thm:cap-identifiability}). This is compatible with the
unrestricted problem, whose cap may increase.

\subsection{Reading guide}

The main distributional proof is in
Sections~\ref{sec:additive}--\ref{sec:matching-lower}.
The supporting material is organized as follows.

\begin{center}
\small
\begin{tabularx}{\textwidth}{@{}p{0.48\textwidth}X@{}}
\toprule
Proof component & Location\\
\midrule
Observable classifiers and multilevel allocation
 & Appendix~\ref{sec:uf-foundations}\\
Weighted sign testing, scale gate, annular saturation
 & Appendix~\ref{sec:lower-foundations}\\
Common hard pairs, adaptive work accounting, additive localization
 & Appendix~\ref{sec:transfer-foundations}\\
Stopping-flow representation and observable realization
 & Appendix~\ref{sec:flow-foundations}\\
Exact robust dual and adaptive comparison
 & Appendices~\ref{sec:exact-dual}--\ref{sec:midpoint-dual}\\
Convergent finite certificates and fixed-cap identification
 & Appendices~\ref{sec:genie}--\ref{sec:identifiability}\\
Exact numerical certificates and reproduction
 & Appendix~\ref{sec:computation}\\
\bottomrule
\end{tabularx}
\end{center}

\section{Related work and the scope of the contribution}
\label{sec:related-work}

\paragraph{Adaptive mixtures of coins.}
Lee--Valiant is the closest direct antecedent \cite{lv}. In the extended
version, Sections~4--5 develop triangular single-coin estimators and
single-coin-to-population lower-bound decompositions. Section~7,
Theorems~4--5 and Lemma~22, optimizes an unbiased estimator when both
sign-conditional bias distributions, a nominal mixture weight, and a
per-coin cap are supplied. Its linear program optimizes information per
expected sample. Section~5.2 also reveals a coin's true bias after a cap
in a strengthened experiment.

Knowing our marginal magnitude law does not supply those conditional
distributions. For $0<\theta(q)<1$, the corresponding magnitude laws are
\[
 \mu_q^+(dd)=\frac{q(d)}{\theta(q)}\,\mu(dd),\qquad
 \mu_q^-(dd)=\frac{1-q(d)}{1-\theta(q)}\,\mu(dd).
\]
They change with the unknown function $q$. Our optimization must use one
observable procedure whose bias, moment, and expected work bounds hold
throughout $\mathcal Q_\rho$. The exact dual therefore retains separate
endpoint, moment-reference, and work-reference populations. The
additional statements concern that robust program, its comparison with
the full population experiment, and convergent upper and lower
certificates. Triangular flows, information/work optimization, adaptive
decomposition, and latent revelation are credited to their antecedents.

\paragraph{Gap entropy and sign testing.}
Chen--Li introduced entropy of normalized dyadic inverse-square gap
costs \cite[Definition~3.3]{cl16}. Chen--Li--Qiao proved the weighted
Sign-$\xi$ entropy inequality \cite[Lemma~4.1]{clq}; Appendix~D develops
its stopping-window argument, and Remark~D.4 explicitly covers Bernoulli
families away from their boundary. Our weighted sign-testing lemma is
an adaptation of this argument. In particular, the form
$A[1+\Ent(\pi)]$ and the statistical origin of its entropy factor are
established ideas.

The additional issue here is how much of a known magnitude law should
pay that entropy cost. The optimizing set in
\eqref{eq:scalar-selection} depends jointly on the law, rarity, and
budget. Its complement pays a truncated inverse-gap tail cost. The
upper construction implements this decomposition without observing
component membership; the lower construction respects both rarity
and mass saturation while combining disjoint components. These steps
are needed to obtain the all-law selection formula from the sign-testing
ingredient. Best-arm identification and latent population-proportion
estimation have different targets and sampling access.

\paragraph{Convex estimation, testing, and stopping.}
Juditsky--Nemirovski give constant-factor affine estimation under their
good-pair assumptions, with a Hellinger-affinity characterization
\cite[Theorem~3.1 and Proposition~3.1]{jn}. Polyanskiy--Wu characterize
quadratic minimax risk for an affine functional over a convex prior
class through a $\chi^2$ modulus for a fixed iid observation kernel
\cite[Theorem~1]{pw}. These results provide the relevant
estimation--testing duality. For a fixed capped one-unit policy, its
transcript is a fixed observation kernel. Optimizing that policy while
controlling expected work uniformly over $q$, and then comparing with
arbitrarily interleaved population algorithms, requires the additional
arguments given here. Strong convex duality and occupation-measure
formulations of optimal stopping are standard tools \cite{bv,cs}; the
four-prior formula is a specific application, not a new general duality
principle.

\paragraph{Baseline and narrower claims.}
The supplied manuscript and pinned version-5 development already
contain the balanced-law benchmark, rare-event phase results, the
variational $\Psi$ lower bound, and the critical quadratic-law entropy
correction \cite{original,repo}. The pinned critical proof also credits
Chen--Li--Qiao for its sign-entropy mechanism. Those results remain baseline contributions. The specialized ingredients
needed here are proved within this manuscript, with their attribution
preserved.

The principal added claim is the joint, arbitrary-law entropy--tail
characterization, together with a robust stopping-program formulation
that admits convergent certificates. The fixed-cap odd-moment
identification result is a secondary consequence of this model and
elementary moment interpolation; it is not presented as a general
mixture-identification principle. The targeted primary-source comparison
above identifies the closest checked statements and their hypotheses.
It does not certify publication priority.

\section{Entropy versus a truncated inverse-gap tail}\label{sec:additive}

This section defines a scalar characterization that separates the part of
the law handled by rare-event scale entropy from a part handled by the
ordinary inverse-gap cost and a part left unresolved at the requested
accuracy. We prove its observable upper bound here and its matching lower
bound in Section~\ref{sec:matching-lower}. The decomposition is analytical:
the algorithm never observes a coin's component label.

\subsection{A homogeneous entropy and a selection functional}

For a finite submeasure $\sigma\le\mu$, use the canonical bands
$I_j=(b_j/2,b_j]$, $b_j=2^{-(j+1)}$, and put
\[
 a_j(\sigma)=\int_{I_j}d^{-2}\,\sigma(dd),\qquad
 A_\sigma=\sum_j a_j(\sigma).
\]
For a finite nonnegative vector $a$, define
\[
 \mathcal G(a)=A+\sum_j a_j\log(A/a_j),\qquad A=\sum_j a_j,
 \qquad \mathcal G(0)=0.
\]
Zero coordinates contribute zero. The functional $\mathcal G$ is concave, coordinatewise increasing,
homogeneous of degree one, and superadditive. Indeed, on positive
coordinates its derivative is $1+\log(A/a_j)$ and its Hessian is
$A^{-1}\mathbf1\mathbf1^\top-\operatorname{diag}(a_j^{-1})$,
which is negative semidefinite by Cauchy--Schwarz. Continuity gives the
boundary cases, and concavity plus homogeneity gives superadditivity. Define
\begin{equation}\label{eq:submeasure-entropy}
 \Gamma_\sigma=\sup_J\mathcal G(a_0(\sigma),\ldots,a_J(\sigma)).
\end{equation}
Thus $\Gamma_0=0$, and, when $0<A_\sigma<\infty$,
$\Gamma_\sigma=A_\sigma[1+\Ent((a_j/A_\sigma)_j)]$, with extended-value
interpretation. Also write
$H_\tau=\int d^{-1}\,\tau(dd)$ for any finite submeasure $\tau$.

Let $h=B^{-1/2}$ and define
\begin{equation}\label{eq:scalar-selection}
 \boxed{\quad
 \mathcal K_{\mu,\rho}(B)=
 \min\left\{\rho,\inf_{0\le\sigma\le\mu}
 \left[h\sqrt{\rho\Gamma_\sigma}
       +\int\min\{1,h/d\}\,d(\mu-\sigma)\right]\right\}.
 \quad}
\end{equation}
Equivalently,
\begin{equation}\label{eq:three-piece-functional}
 \mathcal K_{\mu,\rho}(B)=
 \min\left\{\rho,\inf_{\sigma+\tau+\zeta=\mu}
 \left[\zeta((0,1/2])
       +\frac{\sqrt{\rho\Gamma_\sigma}+H_\tau}{\sqrt B}\right]\right\}.
\end{equation}
Indeed, at fixed $\sigma$, assign each remaining magnitude to the cheaper
of the costs $1$ and $h/d$.

\begin{lemma}[Hard assignments suffice]\label{lem:hard-selection}
The infimum in \eqref{eq:scalar-selection} is unchanged if
$\sigma=\mu|S$ for a measurable set $S$. In particular, for a finite
atomic law a minimizing assignment can put each atom entirely in one of
the three pieces in \eqref{eq:three-piece-functional}.
\end{lemma}

\begin{proof}
Write $f=d\sigma/d\mu\in[0,1]$ and $S_t=\{f>t\}$, so
$\sigma=\int_0^1\mu|S_t\,dt$. Concavity of the square root of every
finite-fragment entropy gives
\[
 \sqrt{\Gamma_\sigma}\ge
       \int_0^1\sqrt{\Gamma_{\mu|S_t}}\,dt
\]
after monotone convergence in the fragment size. For a finite-objective
input the relevant entropies are finite for $t>0$, since
$\mu|S_t\le\sigma/t$ and
$\Gamma_{\mu|S_t}\le\Gamma_\sigma/t$. The remainder integral is linear,
so the objective at $\sigma$ is at least the average of the hard-set
objectives. Their infimum is therefore no larger. The reverse inequality
is immediate. For a finite law the minimum is attained among the finitely
many assignments. This is a selection problem, not a claim of convex
minimization over its assignments.
\end{proof}

The definition immediately implies
\[
 \mathcal K_{\mu,\rho}(B)\le
 \min\{\rho,\omega_\mu(B^{-1/2}),
                  \sqrt{\rho\Gamma_\mu/B}\},
 \qquad
 \omega_\mu(h)=\int\min\{1,h/d\}\,\mu(dd).
\]
Infinite quantities are allowed in this display.

\begin{theorem}[An entropy--tail law for the full joint problem]
\label{thm:entropy-tail-law}
There are universal constants $0<c<C<\infty$ such that, for every
probability measure $\mu$ on $(0,1/2]$, every $0<\rho\le1/4$, and every
$B>0$,
\begin{equation}\label{eq:entropy-tail-law}
 \boxed{\qquad
 c\,\mathcal K_{\mu,\rho}(B)\le R_{\mu,\rho}(B)
             \le C\,\mathcal K_{\mu,\rho}(B).
 \qquad}
\end{equation}
The functional has the equivalent selection and three-piece forms
\eqref{eq:scalar-selection}--\eqref{eq:three-piece-functional}, and hard
measurable assignments suffice. No density, regular variation, finite
inverse moment, or finite full-law entropy is required.
\end{theorem}

The upper half is Theorem~\ref{thm:three-piece-upper}; the lower half is
Theorem~\ref{thm:entropy-tail-lower}. The comparison constants in
\eqref{eq:entropy-tail-law} are universal but conservative. They are not
the factor sixteen proved for the finite convex program; these are
different representations and different comparison bounds.
The proof also gives the explicit conservative constants in
\eqref{eq:explicit-main-constants}.

\begin{theorem}[Three-piece observable upper bound]\label{thm:three-piece-upper}
There is a universal constant $C$ such that, for every magnitude law on
$(0,1/2]$, every $0<\rho\le1/4$, and every $B>0$,
\begin{equation}\label{eq:three-piece-upper}
 R_{\mu,\rho}(B)\le C\mathcal K_{\mu,\rho}(B).
\end{equation}
No component label or magnitude is revealed to the algorithm.
\end{theorem}

\subsection{An uncapped reference bank with a capped observable race}

Fix a decomposition $\mu=\sigma+\tau+\zeta$, put
$z=\zeta((0,1/2])$, and assume $\Gamma_\sigma,H_\tau<\infty$.
It suffices to treat $z\le\rho/128$: otherwise the constant $\rho/2$
has risk at most $64z$. Write $A=A_\sigma$, $H=H_\tau$.
If $\sigma\ne0$, use all its occupied bands, with weights
$\pi_j=a_j(\sigma)/A$ and references $b_j/2$ and $b_j/4$.
For an internal allowance $u<\rho$, set $\epsilon=u/\rho$.
Stop negative if a negative reference likelihood against fair reaches
$1/(\epsilon\pi_j)$; stop positive if
\begin{equation}\label{eq:split-bank-ratio}
 \frac{\sum_j\pi_jP_{1/2+b_j/4,n}}
      {\int P_{1/2-d,n}\,\mu(dd)}\ge1/u.
\end{equation}
Give negative crossings priority. A zero denominator triggers a positive
stop. The split $\sigma+\tau+\zeta=\mu$ specifies the references and the
work analysis; membership in these submeasures is never observed.
The known all-negative mixture in the denominator is the full law
$\mu$, regardless of the decomposition.

By Lemma~\ref{lem:uf-mixture-bank}, a wrong negative crossing has
conditional positive-sign probability at most $\epsilon$, and a positive
crossing has all-negative-mixture probability at most $u$. Thus the total
probability of a wrong bank stop is at most $2u$ over $\mathcal Q_\rho$.
Termination outside $\sigma$ is not assumed.

Race this bank against the ordinary classifier of Lemma~\ref{lem:uf-ordinary}, using
alternating observations of the same physical coin. At an ordinary cap
$M$, its conditional work certificate is
\[
 o_M(d)=\min\{M,17d^{-2}[1+\log_+(Md^2)]\}.
\]
For any finite measure $\xi$, put
\[
 \eta_\xi(M)=\int\min\{1,(Md^2)^{-2}\}\,\xi(dd),\qquad
 O_\xi(M)=\int o_M(d)\,\xi(dd).
\]
The ordinary error under $\mu$ is at most $\eta_\mu(M)$.
Stop as soon as either component stops and return that output.
On completed counterfactual streams, the physical work is at most
\[
 2\min\{T_{\rm bank},T_{\rm ordinary}\}\le2M.
\]
A wrong output implies either a wrong bank stop or an ordinary error.
When $\sigma=0$, omit the bank and use the ordinary classifier alone.

On a $\sigma$-almost-every gap $d$ in band $j$, the matching negative
reference has drift at least $4d^2/3$ and log overshoot below one.
Its expected crossing time is at most
\begin{equation}\label{eq:split-negative-work}
 \frac3{4d^2}[1+\log(1/\epsilon)+\log(1/\pi_j)].
\end{equation}
On a positive gap, the event $\widehat p_n-1/2\ge d/2$ gives
\begin{equation}\label{eq:split-null-laplace}
 \frac{P_{-,\mu,n}}{P_{{\rm f},n}}
 \le\int e^{-2nx^2}\,\mu(dx)
 \le z+\frac A{2en}+\frac H{2\sqrt{en}}.
\end{equation}
The last two bounds use, respectively,
$\sup_xx^2e^{-2nx^2}=1/(2en)$ and
$\sup_xxe^{-2nx^2}=1/(2\sqrt{en})$.
For $n\ge d^{-2}$ the denominator is at most $z+Ad^2+Hd$,
while the matching positive reference contributes at least
$\pi_je^{nd^2/3}$ relative to fair. Equation~\eqref{eq:uf-bank-positive} therefore bounds positive expected bank work by
\begin{equation}\label{eq:split-positive-work}
 \frac6{d^2}\left[
 1+\log(1/\epsilon)+\log(1/\pi_j)
 +\log_+\frac{z+Ad^2+Hd}{\rho}\right].
\end{equation}
These arguments prove finite conditional bank work on $\sigma$.
On $\tau+\zeta$, the ordinary cap guarantees termination of the race.

Since $z/\rho\le1$,
\[
 \log_+\frac{z+Ad^2+Hd}{\rho}
 \le\log3+\log_+(Ad^2/\rho)+\log_+(Hd/\rho).
\]
Using $\sup_{x\ge1}\log x/x=1/e$ and
$\sup_{x\ge1}\log x/x^2=1/(2e)$ gives
\begin{equation}\label{eq:split-surcharge}
 \int q(d)d^{-2}\log_+\frac{z+Ad^2+Hd}{\rho}\,\sigma(dd)
 \le(\log3+1/e)A+\frac{H^2}{2e\rho}.
\end{equation}
Here $\int q\,d\sigma\le\rho$ and $q\le1$ follow from the original
global rarity class; no separate conditional rarity assumption is used.
Integrating \eqref{eq:split-negative-work} and
\eqref{eq:split-positive-work} therefore bounds the bank work on $\sigma$ by
\begin{equation}\label{eq:split-bank-work}
 6\Gamma_\sigma+6A\log(1/\epsilon)
       +6(\log3+1/e)A+\frac{3H^2}{e\rho}.
\end{equation}

\subsection{A resolved-law cap schedule}

Let $\nu=\sigma+\tau$ and set
\[
 \delta_\ell=\rho16^{-(\ell+1)},\qquad
 u_\ell=\delta_\ell/4,\qquad \epsilon_\ell=u_\ell/\rho.
\]
Choose $M_\ell$ as the least power of four, at least four, such that
$\eta_\nu(M_\ell)\le\delta_\ell/8$.
Every chosen cap has a predecessor because $\eta_\nu(4)=1-z$.
The inequality $\eta_\nu(4M)\ge\eta_\nu(M)/16$ implies
\begin{equation}\label{eq:split-cap-selection}
 \delta_\ell/128<\eta_\nu(M_\ell)\le\delta_\ell/8.
\end{equation}
Thus the selected powers of four are distinct and increasing.

If $z>0$, let $J$ be the largest integer with $\delta_J\ge8z$.
Then
\begin{equation}\label{eq:split-terminal-level}
 8z\le\delta_J<128z.
\end{equation}
For $\ell\le J$,
$\eta_\mu(M_\ell)\le\eta_\nu(M_\ell)+z\le\delta_\ell/4=u_\ell$.
The raced classifier has total error at most $3u_\ell\le\delta_\ell$
and false-positive probability at most $2u_\ell\le\delta_\ell$.
When $z=0$, these conclusions hold at every level.

Since $\log(1/\epsilon_\ell)=\log64+\ell\log16$ and
$A\le\Gamma_\sigma$, \eqref{eq:split-bank-work} is at most
$40(\ell+1)\Gamma_\sigma+2H^2/\rho$.
A uniform raced-work certificate is consequently
\begin{equation}\label{eq:split-level-work}
 C_\ell=80(\ell+1)\Gamma_\sigma+4H^2/\rho
                  +2O_\tau(M_\ell)+2zM_\ell.
\end{equation}
It remains a valid, possibly loose, bound when $\sigma=0$ and the bank is
omitted.

The binary level-zero output has second moment at most
$\rho+\delta_0=17\delta_0$. An adjacent difference, run on successive
unused blocks of the same coin, has second moment at most
$\delta_\ell+\delta_{\ell-1}=17\delta_\ell$.
Set $v_\ell=17\delta_\ell$, $k_0=C_0$,
$k_\ell=C_\ell+C_{\ell-1}$ for $\ell\ge1$, and
$S_J=\sum_{\ell\le J}\sqrt{v_\ell k_\ell}$.
Independent-Poisson allocation across these corrections
(Lemma~\ref{lem:uf-poisson})
gives RMS at most
\begin{equation}\label{eq:split-poisson}
 \sqrt{\delta_J^2+S_J^2/B}.
\end{equation}
The seed may choose zero replications, and the Poisson means may be
below one. Every first and subsequent observation is charged.
Every such finite-level procedure has finite deterministic local caps.
The reference bank may be countable; its stopping events are measurable
countable unions and its positive reference is a probability mixture.

\subsection{Controlling mixed work without an extra logarithm}

For every finite $\tau$, the following bounds hold:
\begin{equation}\label{eq:split-elementary-kernels}
 O_\tau(M)\le21H\sqrt M,\qquad
 \eta_\sigma(M)\le A/M,\qquad
 \eta_\tau(M)\le H/\sqrt M.
\end{equation}
The latter two are pointwise inequalities before integration.
For the first, put $x=\sqrt M\,d$. The ratio of $o_M(d)$ to
$\sqrt M/d$ is at most one for $x\le1$, and at most
$17(1+2\log x)/x\le34/\sqrt e<21$ for $x\ge1$.
Lemma~\ref{lem:uf-kernel} supplies the homogeneous ordinary kernel bound:
\begin{equation}\label{eq:split-ordinary-kernel}
 \sum_{k\ge1}\sqrt{\eta_\tau(4^k)O_\tau(4^k)}\le 208H.
\end{equation}
Its proof applies without normalization of $\tau$: dyadic binning yields
the summable kernels $2^{k-j}$ for $j\ge k$, and
$2^{3(j-k)}$, $2^{j-k}(1+k-j)$ for $j<k$.
Apply $\sqrt{xy}\le(x+y)/2$ and sum these kernels.

Split levels according to which term of
$\eta_\nu=\eta_\sigma+\eta_\tau$ is at least half the total.
On levels where $\eta_\tau$ dominates,
$\delta_\ell\le256\eta_\tau(M_\ell)$, so distinctness of the caps and
\eqref{eq:split-ordinary-kernel} give a total
$\sum\sqrt{\delta_\ell O_\tau(M_\ell)}\le3328H$.
On the other levels,
\[
 \delta_\ell\le256\eta_\sigma(M_\ell)\le256A/M_\ell,
 \qquad M_\ell\le256A/\delta_\ell.
\]
Therefore
\[
 \sqrt{\delta_\ell O_\tau(M_\ell)}
 \le4\sqrt{21}\,H^{1/2}A^{1/4}\delta_\ell^{1/4}.
\]
The sum of $\delta_\ell^{1/4}$ over all levels is $\rho^{1/4}$.
The arithmetic--geometric mean inequality proves
\begin{equation}\label{eq:split-mixed-sum}
 \sum_{\ell\le J}\sqrt{\delta_\ell O_\tau(M_\ell)}
 \le C_{\rm mix}[H+\sqrt{\rho A}].
\end{equation}

The work on the unresolved mass is controlled separately.
On levels where $\eta_\sigma$ dominates,
\begin{equation}\label{eq:split-unresolved-sigma}
 \sum\sqrt{z\delta_\ell M_\ell}
 \le16(J+1)\sqrt{zA}\le\sqrt{2\rho A}.
\end{equation}
The last inequality uses $z\le\delta_J/8$ and
$(J+1)4^{-(J+1)}\le1/4$.
On levels where $\eta_\tau$ dominates,
$\sqrt{M_\ell}\le256H/\delta_\ell$, so
\begin{equation}\label{eq:split-unresolved-tau}
 \sum\sqrt{z\delta_\ell M_\ell}
 \le256H\sqrt z\sum_{\ell=0}^J\delta_\ell^{-1/2}
 \le\frac{1024}{3}H\sqrt{z/\delta_J}
 \le\frac{1024}{3\sqrt8}H.
\end{equation}
When $z=0$ these contributions vanish.

\begin{proof}[Completion of Theorem~\ref{thm:three-piece-upper}]
The shifted adjacent-level sum has a factor $1/4$, so
\[
 S_J\le\frac{5\sqrt{17}}4
                \sum_{\ell=0}^J\sqrt{\delta_\ell C_\ell}.
\]
Use \eqref{eq:split-level-work}, \eqref{eq:split-mixed-sum},
\eqref{eq:split-unresolved-sigma}, \eqref{eq:split-unresolved-tau},
and $A\le\Gamma_\sigma$. The elementary sums
$\sum_{\ell\ge0}4^{-\ell}\sqrt{\ell+1}\le16/9$ and
$\sum_{\ell\ge0}4^{-\ell}=4/3$ then give
\[
 S_J\le C_{\rm sum}[\sqrt{\rho\Gamma_\sigma}+H].
\]
For $z>0$, \eqref{eq:split-terminal-level} and
\eqref{eq:split-poisson} imply
$R(B)\le128z+C_{\rm sum}[\sqrt{\rho\Gamma_\sigma}+H]/\sqrt B$.
For $z=0$, the bound on $S_J$ holds at every finite terminal level.
Use \eqref{eq:split-poisson} and then $\delta_J\to0$; taking the infimum
over these finite constructions gives the asserted bound.
The constant estimator covers $z>\rho/128$.
Finally take the infimum over decompositions and compare with $\rho/2$.
All approximations are at fixed $(\mu,\rho,B)$. The comparison constant
is independent of the chosen decomposition and terminal level; no
attaining infinite-dimensional decomposition is required.
No infinite-algorithm interchange is required. The same infimum can be
approached with finite reference banks: finite $\Gamma_\sigma$ implies
$H_\sigma<\infty$, so truncate $\sigma$ to finitely many bands and move its
remainder to $\tau$. The entropy does not increase and the additional
inverse-gap cost tends to zero. Then use the finite-level construction.
\end{proof}

\section{A matching scalar lower bound by entropy packing}
\label{sec:matching-lower}

The lower proof uses optimality of the retained set in
\eqref{eq:scalar-selection}. Its ordinary and ignored pieces are controlled
by the variational modulus proved in Appendix~\ref{sec:transfer-foundations}. Its retained entropy can be packed
into disjoint components on which the actual annular mass constraint is
active. The additive localization theorem then combines these components
at every budget.

\begin{theorem}[Matching entropy--tail lower bound]
\label{thm:entropy-tail-lower}
There is a universal $c>0$ such that, for every magnitude law,
$0<\rho\le1/4$, and $B>0$,
\[
 R_{\mu,\rho}(B)\ge c\,\mathcal K_{\mu,\rho}(B).
\]
Together with Theorem~\ref{thm:three-piece-upper}, this proves
Theorem~\ref{thm:entropy-tail-law}.
\end{theorem}

\subsection{Canonical endpoints change the functional by at most two}

Put $h=B^{-1/2}$, $w_j=\mu(I_j)$, and $a_j=w_j/b_j^2$, where
$b_j=2^{-(j+1)}$ and $I_j=(b_j/2,b_j]$.
For a finite prefix, define
\begin{equation}\label{eq:canonical-selection-prefix}
 k_J=\min\left\{\rho,\min_{I\subset\{0,\ldots,J\}}
 \left[h\sqrt{\rho\,\mathcal G((a_j)_{j\in I})}
       +\sum_{\substack{j\le J\\j\notin I}}
                    w_j\min\{1,h/b_j\}\right]\right\}.
\end{equation}
Fractional retained band masses cannot improve this finite expression:
the square root of $\mathcal G$ is concave, the remainder is linear,
and the minimum over a finite splitting box occurs at a vertex.
The values $k_J$ increase to a limit $k$, with
\begin{equation}\label{eq:canonical-prefix-limit}
 k_J\le k\le\min\{\rho,k_J+\sum_{j>J}w_j\}.
\end{equation}
Restriction proves the first inequality; assigning every omitted band to
ignored mass proves the second. The limit is at fixed $(\mu,\rho,B)$;
no law-independent truncation depth or nested choice of minimizers is
required. The omitted tail probability must be retained before the limit.

The exact and canonical functionals satisfy
\begin{equation}\label{eq:canonical-factor-two}
 k\le\mathcal K_{\mu,\rho}(B)\le2k.
\end{equation}
Indeed, on a band the exact inverse-square cost lies between its
canonical cost and four times that cost, while
\[
 \min\{1,h/b_j\}\le\min\{1,h/d\}\le2\min\{1,h/b_j\}.
\]
Monotonicity and homogeneity of $\mathcal G$ give the factor two after
taking the square root. The lower comparison applies to any fractional
submeasure before minimizing the canonical splitting box. For the upper
comparison, lift a canonical choice to a union of whole original bands
and assign the vanishing omitted tail to ignored mass.
These are comparisons of functionals. We do not replace the original
observation experiment by an experiment with rounded magnitudes.

We prove $R_{\mu,\rho}(B)\ge c k_J$ uniformly in the finite prefix.
The limit \eqref{eq:canonical-prefix-limit} and
\eqref{eq:canonical-factor-two} will complete the theorem.

\subsection{Two consequences of optimal finite selection}

Fix a minimizer $I$ of the inner finite objective in
\eqref{eq:canonical-selection-prefix}, before its outer cap by $\rho$.
Let $O$ consist of the complementary bands with $b_j\ge h$, and $Z$ of
those with $b_j<h$. Write
\[
 F=\mathcal G((a_j)_{j\in I}),\qquad A=\sum_{j\in I}a_j,\qquad
 z=\sum_{j\in Z}w_j,\qquad H_O=\sum_{j\in O}a_jb_j.
\]
The minimized inner cost is $z+hH_O+h\sqrt{\rho F}$.
Zero-mass bands can be omitted.

\begin{lemma}[Optimal-band inequalities]\label{lem:optimal-band-inequalities}
Every ordinary band $j\in O$ has $w_j\le\rho$.
If $F>0$ and $\lambda=\sqrt{\rho/F}$, every retained band satisfies
\begin{equation}\label{eq:retained-gap-condition}
 b_j\ge\frac{\lambda}{2}[1+\log(A/a_j)].
\end{equation}
\end{lemma}

\begin{proof}
For an ordinary band, put $x=a_j$ and compare with moving it to $I$.
Optimality gives
\[
 xb_j\le\sqrt\rho[\sqrt{F+\Delta}-\sqrt F].
\]
If $F=0$, the squared bracket is $x$. Otherwise use
$\log(1+t)\le\sqrt t$ for $t\ge0$ to obtain
\[
 \Delta=x+A\log(1+x/A)+x\log(1+A/x)
       \le x+2\sqrt{Ax}\le x+2\sqrt{Fx}.
\]
Thus $\sqrt{F+\Delta}-\sqrt F\le\sqrt x$, and optimality gives
$x^2b_j^2\le\rho x$, proving $w_j\le\rho$.

For $j\in I$, removing the band and assigning it to ordinary work is
admissible even if ignoring it would be cheaper. Hence
\[
 a_jb_j\ge\sqrt\rho[\sqrt F-\sqrt{F_{-j}}].
\]
The derivative of $\mathcal G$ in this coordinate is
$1+\log(A/a_j)$; it only increases as that coordinate is removed.
Therefore
$F-F_{-j}\ge a_j[1+\log(A/a_j)]$.
Dividing by $\sqrt F+\sqrt{F_{-j}}\le2\sqrt F$ proves
\eqref{eq:retained-gap-condition}.
\end{proof}

\subsection{The ordinary and ignored costs are already lower bounds}

Recall the functional from Proposition~\ref{prop:psi-null}:
\[
 \Psi_{\mu,\rho}(h)=
 \sup_{\substack{0\le q\le1\\\int q\,d\mu\le\rho}}
 \int\min\{q(d),h\sqrt{q(d)}/d\}\,\mu(dd).
\]
Theorem~\ref{thm:bias-transfer} and Proposition~\ref{prop:psi-null} imply
\begin{equation}\label{eq:psi-budget-lower}
 R_{\mu,\rho}(B)\ge\Psi_{\mu,\rho}(B^{-1/2})/8192.
\end{equation}
For an explicit deduction, take $b=\Psi(h)/64\le\rho/64$ in the
bias-indexed null obstruction in Proposition~\ref{prop:psi-null}. Evaluating it at
$h=B^{-1/2}$ gives
$J\ge(15\Psi(h)/16)^2/(192h^2)$.
The bias-indexed transfer has its smaller branch at $b$, yielding
\eqref{eq:psi-budget-lower}.

Every ignored band has $d\le b_j<h$, so allocating on these bands gives
\begin{equation}\label{eq:ignored-psi}
 \min\{\rho,z\}\le\Psi_{\mu,\rho}(h).
\end{equation}
For the ordinary bands we need an elementary dyadic inequality.
If masses $v_j$ satisfy $v_j\le s$ at values
$f_j=b_j^{-1}=2^{j+1}$, then
\begin{equation}\label{eq:dyadic-top-mass}
 \sum_jv_jf_j\le2
 \sup_{\substack{0\le q_j\le1\\\sum_jv_jq_j\le s}}
                 \sum_jv_jq_jf_j.
\end{equation}
If the total mass is at most $s$, this is immediate. Otherwise select
the largest $f_j$ up to mass $s$, allowing a fraction at the boundary
index $j_0$. Let $r$ be the selected mass at the boundary index.
The selected mass above it is $s-r$, at values at least $2f_{j_0}$,
so the selected integral is at least $(2s-r)f_{j_0}$.
The excluded smaller values contribute at most $sf_{j_0}$, and the
excluded boundary mass contributes at most $(s-r)f_{j_0}$.
Thus the excluded integral is no larger than the selected integral,
proving \eqref{eq:dyadic-top-mass}.

Apply it to $O$ with $s=\rho$, using
Lemma~\ref{lem:optimal-band-inequalities}. It follows that $H_O$ is at
most twice the largest ordinary-band inverse-gap integral at mass $\rho$.
On these bands $h/b_j\le1$ and $d\le b_j$;
for every $0\le q\le1$,
$\min\{q,h\sqrt q/d\}\ge qh/b_j$.
Consequently
\begin{equation}\label{eq:ordinary-psi}
 hH_O\le2\Psi_{\mu,\rho}(h).
\end{equation}

\subsection{Packing entropy into comparable-cost groups}

\begin{lemma}[Entropy tiers]\label{lem:entropy-tiers}
For a finite positive vector $a$, let $M=\max_j a_j$ and group its
coordinates by $a_j\in(M2^{-k-1},M2^{-k}]$, $k\ge0$.
If $F_k$ is each group's entropy cost, then
\begin{equation}\label{eq:entropy-tiers}
 \mathcal G(a)\le7\sum_k F_k.
\end{equation}
\end{lemma}

\begin{proof}
Let $A_k$ be the cost mass in group $k$, $n_k$ its cardinality, and
$p_k=A_k/A$. Entropy decomposition gives
$\mathcal G(a)=\sum_kF_k+A\Ent(p)$.
Comparison with the probability vector proportional to $2^{-k/2}$ gives
\[
 A\Ent(p)\le A\log(2+\sqrt2)
                 +\frac{\log2}{2}\sum_k kA_k.
\]
On groups with $n_k\ge2^{k/2}$, the second term is at most
$\sum_kA_k\log n_k\le\sum_kF_k$; comparable weights imply
$F_k\ge A_k[1+\log n_k-\log2]\ge A_k\log n_k$.
On the other groups, $A_k\le M2^{-k/2}$, so their contribution is at most
\[
 \frac{\log2}{2}M\sum_{k\ge0}k2^{-k/2}<3M\le3A.
 \]
Finally $\log(2+\sqrt2)<2$ and $A\le\sum_kF_k$.
Adding the initial $\sum_kF_k$ proves \eqref{eq:entropy-tiers}.
\end{proof}

Split the retained low-gap bands $j\ge6$ by index parity and then by
these cost tiers. Parity splitting loses at most a factor two in
entropy, so all resulting group entropies sum to at least $F_{\rm low}/14$.
Every group is a separated annular family with endpoint ratios at least
four, and its cost weights lie in some interval $(c,2c]$.

Consider a group of $n$ bands and cost entropy at most $2cn(1+\log n)$.
Since $A\ge nc$ and every group coordinate is at most $2c$,
\eqref{eq:retained-gap-condition} implies
\[
 b_j\ge\frac{\lambda}{2}[1+\log(n/2)]
       \ge\frac{\lambda}{8}(1+\log n).
\]
Remove the $t=\lceil\log_4n\rceil$ smallest-gap bands, with $t=0$ for
$n=1$. For $n\ge2$, $t\le n/2$. The remaining entropy $G_i$ satisfies
\begin{equation}\label{eq:trimmed-entropy}
 G_i\ge F_{\rm group}/16.
\end{equation}
Indeed, at least $n/2$ coordinates remain, each at least $c$.
Coordinate monotonicity gives entropy at least
$cn'(1+\log n')$, whereas the original entropy is at most
$2cn(1+\log n)$. Since
$1+\log(n/2)\ge(1-\log2)(1+\log n)$, their ratio is at least
$(1-\log2)/4>1/16$. The case $n=1$ is immediate.

Separation amplifies the smallest remaining endpoint by at least
$4^t\ge n$. Hence
\begin{equation}\label{eq:trimmed-minimum-gap}
 b_{\min,i}\ge\lambda n(1+\log n)/8.
\end{equation}
Here $G_i$ is the canonical endpoint entropy of the retained bands,
which is exactly the annular quantity $\Gamma_{\sigma_i}^{\circ}$
in Appendix~\ref{sec:lower-foundations}.
Let $\sigma_i$ be the original law restricted to these remaining bands,
$\tau_i$ its mass, $a_i$ its largest endpoint, and $w_{*,i}$ its smallest
unnormalized band mass. Then
$\tau_i/a_i^2\ge c$ and $w_{*,i}\ge c b_{\min,i}^2$.
Together with $G_i\le2cn(1+\log n)$, this gives the actual annular
saturation bound
\begin{equation}\label{eq:packed-saturation}
 \boxed{\quad
 \frac{\tau_iw_{*,i}}{a_i^2}
        \ge\frac{\lambda^2G_i^2}{256}.
 \quad}
\end{equation}
This step retains the minimum-band-mass constraint; it does not replace
that constraint by an inverse-gap moment.

Allocate original rarity and conditional rarity by
\begin{equation}\label{eq:packed-rarities}
 \theta_i=\frac{\lambda^2G_i}{1024},
 \qquad r_i=\theta_i/\tau_i.
\end{equation}
Equation \eqref{eq:packed-saturation} makes $\theta_iG_i$ at most one
quarter of the saturation term. Since $G_i\ge\tau_i/a_i^2$, that
saturation divided by $G_i$ is at most $w_{*,i}\le\tau_i$.
Thus $\theta_i\le\tau_i/4$ and $r_i\le1/4$.
The components are disjoint, and superadditivity gives
\[
 \sum_i\theta_i\le\lambda^2F/1024=\rho/1024.
\]
Every retained component is supported on $(0,1/128]$; parity gives
endpoint ratios at least four. The normalized minimum band mass is
$w_{*,i}/\tau_i$, so \eqref{eq:packed-saturation} is exactly the
unnormalized saturation condition required by the annular proposition.
Every component is therefore admissible, and every entropy branch of the
scale-uniform annular lower bound is active.

\subsection{Combining the components at every budget}

Let $F_{\rm high}$ be the retained entropy in the first six bands and
$F_{\rm low}$ the remainder. Entropy decomposition gives
$F\le2(F_{\rm high}+F_{\rm low})$, and
$F_{\rm high}\le4096(1+\log6)$.
If $F_{\rm high}\ge F_{\rm low}$, then $F$ is bounded by an absolute
constant. The feasible constant allocation $q=\rho$ gives
$\Psi(h)\ge\min\{\rho,2h\sqrt\rho\}$, so
\begin{equation}\label{eq:high-entropy-psi}
 \min\{\rho,h\sqrt{\rho F}\}\le C\Psi(h).
\end{equation}

Otherwise $F_{\rm low}>F_{\rm high}$ and $F\le4F_{\rm low}$.
Lemma~\ref{lem:entropy-tiers}, parity splitting, and
\eqref{eq:trimmed-entropy} give
\begin{equation}\label{eq:packed-entropy-total}
 \sum_iG_i\ge F_{\rm low}/224\ge F/896.
\end{equation}
Proposition~\ref{prop:scaled-annular} and
Theorem~\ref{thm:additive-localization} now imply
\[
 R_{\mu,\rho}(B)\gtrsim
       \sum_i\min\{\theta_i,\sqrt{\theta_iG_i/B}\}.
\]
All conditional bias and annular constants here are fixed universal
constants. Crucially, the ratio $\theta_i/G_i=\lambda^2/1024$ is common
to every component. This remains true with the fixed conditional-bias
and annular constants retained, so all components have the same budget
crossover. Therefore the displayed sum equals
\[
 \left(\sum_iG_i\right)
       \min\{\lambda^2/1024,\lambda/(32\sqrt B)\}.
\]
Using \eqref{eq:packed-entropy-total} and
$\lambda^2=\rho/F$ proves
\begin{equation}\label{eq:retained-entropy-lower}
 R_{\mu,\rho}(B)\ge c
                  \min\{\rho,h\sqrt{\rho F}\}.
\end{equation}
No limit in $B$ or separate lower threshold for each component is used.

\begin{proof}[Completion of Theorem~\ref{thm:entropy-tail-lower}]
If $z\ge\rho$, \eqref{eq:ignored-psi} and
\eqref{eq:psi-budget-lower} already give $R\ge c k_J$.
Otherwise \eqref{eq:ignored-psi} and \eqref{eq:ordinary-psi} imply
\[
 k_J\le\min\{\rho,3\Psi(h)+h\sqrt{\rho F}\}
       \le3\Psi(h)+\min\{\rho,h\sqrt{\rho F}\}.
\]
Apply \eqref{eq:psi-budget-lower} and either
\eqref{eq:high-entropy-psi} or \eqref{eq:retained-entropy-lower}.
If $F=0$, the $\Psi$ bound alone suffices.
Thus $R_{\mu,\rho}(B)\ge c k_J$ with one universal constant for every
finite prefix. Let $J$ increase and use
\eqref{eq:canonical-prefix-limit}--\eqref{eq:canonical-factor-two}.
This proves the theorem. The passage to every magnitude law uses only
vanishing omitted tail probability and finite-band lower constructions.
\end{proof}


\subsection{Explicit conservative comparison constants}

The preceding arguments provide, without optimizing constants,
\begin{equation}\label{eq:explicit-main-constants}
 10^{-17}\mathcal K_{\mu,\rho}(B)
 \le R_{\mu,\rho}(B)\le33000\mathcal K_{\mu,\rho}(B).
\end{equation}
For completeness, the work summation in Section~\ref{sec:additive} has
entropy and ordinary coefficients at most, respectively,
\[
 \frac{5\sqrt{17}}4
 \left(\frac{4\sqrt{80}}9+2\sqrt{42}+2\right)<104,
\]
\[
 \frac{5\sqrt{17}}4
 \left(\frac23+\sqrt2(3328+2\sqrt{21})+\frac{512}3\right)<27000.
\]
Together with the $128z$ bias term and the constant estimator, these
give the rounded upper constant.

For the lower bound let $c_{\rm ann}=e^{-28}/2000000$ and
$c_E=\sqrt{c_{\rm ann}}/(896\cdot128\cdot32)$.
Before suppressing constants, the packed contribution is at least
\[
 \frac{\sum_iG_i}{128}
 \min\left\{\frac{\lambda^2}{65536},
             \frac{\sqrt{c_{\rm ann}}\lambda}{32\sqrt B}\right\}
 \ge c_E\min\{\rho,h\sqrt{\rho F}\}.
\]
The high-gap alternative gives
$\min\{\rho,h\sqrt{\rho F}\}\le112\Psi(h)$ and has a stronger
constant. Combining \eqref{eq:psi-budget-lower} with the final
$3\Psi(h)$ assembly and $\mathcal K\le2k$ gives the coefficient
$[2(24576+c_E^{-1})]^{-1}>10^{-17}$.
For example, $e<3$ and $\sqrt{2000000}<1500$ imply
$2(24576+896\cdot128\cdot32\cdot1500\cdot3^{14})<10^{17}$.
These are constants for the scalar comparison, not the factor sixteen
of the stopping-program certificate.

\section{Consequences and limits of the characterization}\label{sec:discussion}

The selection formula recovers two immediate upper bounds by choosing
the retained set to be all or none of the magnitude space:
\[
 R_{\mu,\rho}(B)\lesssim
 \min\left\{\rho,\sqrt{\rho\Gamma_\mu/B},
 \int\min\{1,(d\sqrt B)^{-1}\}\,\mu(dd)\right\}.
\]
The full formula additionally allows different regions of the same law
to use different treatments. Its lower proof explains why contributions
from several irregular components may need to be combined.

For a single magnitude $\mu=\delta_d$, hard assignments leave two choices
and give
\[
 R_{\delta_d,\rho}(B)\asymp
       \min\{\rho,\sqrt{\rho/(Bd^2)}\}.
\]
For a general law, the finite-prefix selection values $k_J$ in
\eqref{eq:canonical-selection-prefix} have an explicit tail-mass error
bound \eqref{eq:canonical-prefix-limit}, and their limit lies between
$\mathcal K/2$ and $\mathcal K$. This is a finite approximation using only
known band masses. Exhaustive selection is possible, but an efficient
general selection algorithm has not been established here.

The scalar characterization concerns the full minimax risk, whereas
$P_N$ bounds worst bias and worst variance separately. Accordingly, its
exact dual is an identity for a certificate program, and the statistical
connection is a comparison inequality. The convergent genie brackets
can certify that a chosen cap suffices for a particular input without
using the large universal cap. They do not give a uniform runtime bound.

The remaining quantitative questions include sharp comparison constants,
efficient selection for structured laws, and rates of cap convergence.
The present theorem assumes known $\mu$, Bernoulli observations, and a
uniform expected observation budget. Extensions to an unknown magnitude
law, other observation kernels, or almost-sure budgets require separate
arguments.

This consolidated manuscript includes the specialized proof chain rather
than relying on the earlier working notes. Separate model audits found
no remaining substantive inconsistency in the reviewed chain; exact
arithmetic verifies the finite numerical certificates. These checks do
not constitute external refereeing or proof-assistant verification of the
general theorems. The related-work comparison is targeted and does not
certify publication priority.


\subsection{Two scope examples}

\paragraph{A noninterval optimal assignment.}
Let $\varepsilon=1/1024$, put masses $(1-\varepsilon)/2$, $\varepsilon$,
and $(1-\varepsilon)/2$ at $1/2$, $1/4$, and $1/8$, respectively, and
take $\rho=1/4$, $h=B^{-1/2}=1/1024$.
All three atoms lie above $h$. Enumerating the eight hard retained sets
shows that the unique minimum retains $\{1/2,1/8\}$ and treats $1/4$
ordinarily. The following rational enclosures for the inner cost divided
by $h$ make the strict comparison explicit:
\begin{center}
\begin{tabular}{@{}lc@{}}
\toprule
Retained atoms & Enclosure for cost divided by $h$\\
\midrule
$\varnothing$ & $(4.99902,4.99903)$\\
$\{1/2\}$ & $(4.70676,4.70677)$\\
$\{1/4\}$ & $(5.05761,5.05762)$\\
$\{1/2,1/4\}$ & $(4.72155,4.72156)$\\
$\{1/8\}$ & $(3.82997,3.82998)$\\
$\{1/2,1/8\}$ & $(3.22748,3.22749)$\\
$\{1/4,1/8\}$ & $(3.83271,3.83272)$\\
$\{1/2,1/4,1/8\}$ & $(3.22943,3.22944)$\\
\bottomrule
\end{tabular}
\end{center}
These follow by substituting the rational weights into $\mathcal G$.
For example, after reducing each logarithm argument to $1\le x\le2$,
use $y=(x-1)/(x+1)$ and
\[
 0\le\log x-2\sum_{k=0}^{m-1}\frac{y^{2k+1}}{2k+1}
 \le\frac{2y^{2m+1}}{(2m+1)(1-y^2)},
\]
together with rational bounds for $\log2$ from the same series and
square-root bounds checked by squaring. Already $m=90$ gives intervals
much narrower than those displayed. The companion verifier
\texttt{checks/verify\_revisions.py} performs these exact calculations.
Since the outer cap is inactive, the original functional has the same
optimizer. No interval can contain both retained atoms while excluding
the intervening positive-mass atom.

\paragraph{Infinite inverse moments and a slow tail.}
Let $\mu\{2^{-j}\}=1/[j(j+1)]$ for $j\ge1$. Both inverse moments
of orders one and two are infinite. At $B=4^J$, ignore $j>J$ and
treat the first $J$ atoms ordinarily. This gives
\[
 \mathcal K(B)\le\frac1{J+1}
       +\sum_{j=1}^J\frac{2^{j-J}}{j(j+1)}
       =\frac1{J+1}+O(J^{-2}).
\]
For fixed $\rho$ and sufficiently large $J$, the ignored-tail allocation
gives $\Psi(2^{-J})\ge1/(J+1)$. Thus the upper and lower scales match
even though the full-law inverse moments diverge. This example uses
only the vanishing tail probability, not an integrable full-law cost.

\clearpage
\appendix
\section{Observable classification and finite Poisson allocation}
\label{sec:uf-foundations}

This section proves the specialized upper-bound tools used below. The
ordinary classifier is the construction in the supplied manuscript
\cite{original,repo}; the weighted reference bank and multilevel
allocation are restated with their proofs so that the present paper
does not depend on the preceding unpublished notes. All likelihoods
refer to the actual observed binary prefix. A reference distribution
is used to evaluate that prefix, and never reveals or redraw the
latent state of a unit.

\subsection{A capped classifier that does not know the gap}

For a dyadic integer $M\ge1$, inspect prefix lengths
$1,2,4,\ldots,M$. At a prefix $n<M$, stop if
\[
 \left|\widehat p_n-\frac12\right|>
 t_{n,M}:=\sqrt{\frac{1+\log(M/n)}{n}},
\]
and output the side of $1/2$ containing $\widehat p_n$. At $M$, output
$\1_{\{\widehat p_M>1/2\}}$, with a negative output on ties.
Denote the output and work by $(Z_M^{\rm ord},T_M^{\rm ord})$.

\begin{lemma}[Ordinary conditional error and work]\label{lem:uf-ordinary}
For every $0<d\le1/2$ and either latent sign,
\begin{align}
 \Pp_{d,s}(Z_M^{\rm ord}\ne s)
 &\le \min\{1,(Md^2)^{-2}\},\label{eq:uf-ordinary-error}\\
 \E_{d,s}T_M^{\rm ord}
 &\le o_M(d):=\min\{M,17d^{-2}[1+\log_+(Md^2)]\}.
 \label{eq:uf-ordinary-work}
\end{align}
\end{lemma}

\begin{proof}
For independent Bernoulli observations of mean $p$, Hoeffding's
inequality gives
$\Pp(\widehat p_n-p\ge t)\le e^{-2nt^2}$ and the analogous lower-tail
bound \cite{hoeffding}. A wrong early output at a fixed inspected prefix
requires a deviation of at least $d+t_{n,M}$. Applying a union bound on
the underlying full tape, without conditioning on the stopping event,
therefore gives
\[
 \Pp_{d,s}(Z_M^{\rm ord}\ne s)
 \le e^{-2Md^2}
       +e^{-2}\sum_{\substack{n<M\\n\ {\rm dyadic}}}
                    (n/M)^2e^{-2nd^2}.
\]
For any dyadic sequence of positive $u$ values, the intervals
$[u/2,u]$ have disjoint interiors and
\[
 \int_{u/2}^{u}y e^{-2y}\,dy
 \ge\frac38u^2e^{-2u}.
\]
Consequently $\sum_u u^2e^{-2u}\le2/3$.
Writing $x=Md^2\ge1$ and using $x^2e^{-2x}\le e^{-2}$, the preceding
error bound is at most $(5/3)e^{-2}x^{-2}<x^{-2}$. For $x<1$, the
probability bound by one suffices.

Since the stopping time takes dyadic values,
\[
 \E_{d,s}T_M^{\rm ord}
 =1+\sum_{\substack{n<M\\n\ {\rm dyadic}}}
                      n\,\Pp_{d,s}(T_M^{\rm ord}>n).
\]
If $d>2t_{n,M}$, survival implies a deviation from the true mean of
at least $d/2$. Thus
\[
 \E_{d,s}T_M^{\rm ord}
 \le1+\sum_{\substack{n<M\\n\ {\rm dyadic}}}
       n\left[\1_{\{d\le2t_{n,M}\}}+2e^{-nd^2/2}\right].
\]
For $x=Md^2\ge1$, the first indicator implies
$nd^2\le4(1+\log x)$. Indeed, this is immediate if $nd^2<1$; otherwise
$\log(M/n)=\log(x/(nd^2))\le\log x$.
The corresponding dyadic sum is at most
$8d^{-2}(1+\log x)$.
Also
\[
 \int_{u/2}^{u}e^{-y/2}\,dy\ge\frac12u e^{-u/2}
 \quad\Longrightarrow\quad
 \sum_u u e^{-u/2}\le4.
\]
The exponential part of the work bound is therefore at most $8d^{-2}$.
Since $1\le d^{-2}$, total expected work is at most
$17d^{-2}(1+\log x)$. Finally $T_M^{\rm ord}\le M$ deterministically;
when $x<1$, the displayed minimum equals $M$.
\end{proof}

For any finite measure $\xi$ on $(0,1/2]$, write
\[
 \eta_\xi(M)=\int\min\{1,(Md^2)^{-2}\}\,\xi(dd),
 \qquad
 O_\xi(M)=\int o_M(d)\,\xi(dd).
\]
Both are homogeneous in $\xi$. The first is nonincreasing in $M$ and
the second is nondecreasing. Pointwise comparison also gives
$\eta_\xi(4M)\ge\eta_\xi(M)/16$.

\begin{lemma}[A homogeneous dyadic kernel estimate]\label{lem:uf-kernel}
For every finite measure $\xi$ and integer $K\ge1$,
\[
 \sum_{k=1}^{K}\sqrt{\eta_\xi(4^k)O_\xi(4^k)}
 \le208\int\min\{2^K,d^{-1}\}\,\xi(dd).
\]
In particular,
\begin{equation}\label{eq:uf-kernel-infinite}
 \sum_{k\ge1}\sqrt{\eta_\xi(4^k)O_\xi(4^k)}
 \le208\int d^{-1}\,\xi(dd),
\end{equation}
with extended-value interpretation.
\end{lemma}

\begin{proof}
Use bands $(2^{-j-1},2^{-j}]$, $j\ge1$, with masses $w_j$, and put
$c_j=2^jw_j$. Direct substitution into the two kernels yields
\[
 \eta_\xi(4^k)\le2^{-k}
 \left[\sum_{j\ge k}c_j2^{k-j}
              +16\sum_{j<k}c_j2^{3(j-k)}\right]=2^{-k}A_k,
\]
\[
 O_\xi(4^k)\le2^k
 \left[\sum_{j\ge k}c_j2^{k-j}
        +136\sum_{j<k}c_j2^{j-k}(1+k-j)\right]=2^kD_k.
\]
For the second display on $j<k$, use
$d^{-2}\le4\,2^{2j}$ and
$1+\log(4^kd^2)\le2(1+k-j)$.
On $j\ge k$, the deterministic cap gives the first sum exactly as
an upper bound.

Apply $\sqrt{A_kD_k}\le(A_k+D_k)/2$ and sum over $k\le K$.
For a fixed $j\le K$, the three relevant geometric sums obey
\[
 \sum_{k\le j}2^{k-j}\le2,\qquad
 \sum_{r\ge1}2^{-3r}=\frac17,\qquad
 \sum_{r\ge1}2^{-r}(1+r)=3.
\]
The coefficient of $c_j$ is thus at most
$2+8/7+204<208$. For $j>K$, only the first kernels occur, with
coefficient at most $2\,2^{K-j}$. Hence the result is at most
\[
 208\sum_{j\le K}2^jw_j+2\,2^K\sum_{j>K}w_j.
\]
For $j\le K$, $\min\{2^K,d^{-1}\}\ge2^j$ on the band; for $j>K$
it equals $2^K$. This proves the finite bound. Increasing $K$ and
using monotone convergence gives \eqref{eq:uf-kernel-infinite}.
All rearrangements involve nonnegative terms, so normalization and
moment assumptions are unnecessary.
\end{proof}

\subsection{A weighted bank compared with the full null mixture}

For a binary prefix $x_{1:n}$, let $P_{p,n}(x_{1:n})$ be its probability
under iid $\operatorname{Ber}(p)$ observations, and let
$P_{{\rm f},n}=P_{1/2,n}$. Choose any nonempty finite or countable set
of canonical bands
\[
 I_j=(b_j/2,b_j],\qquad b_j=2^{-j-1},
\]
and positive weights $\pi_j$ summing to one on these bands.
Set $r_j^-=b_j/2$ and $r_j^+=b_j/4$, and define
\[
 L_{j,n}^-=\frac{P_{1/2-r_j^-,n}}{P_{{\rm f},n}},
 \qquad
 M_n=\frac{\sum_j\pi_jP_{1/2+r_j^+,n}}
           {P_{-,\mu,n}},
 \qquad
 P_{-,\mu,n}=\int P_{1/2-d,n}\,\mu(dd).
\]
For $0<\epsilon<1$, stop negative when any
$L_{j,n}^-\ge1/(\epsilon\pi_j)$, and stop positive when
$M_n\ge1/(\rho\epsilon)$. Give negative crossings priority.
A zero denominator triggers a positive stop. Denote the bank stopping
time by $T_{\rm bank}$; it may be infinite outside its selected bands.

\begin{lemma}[Full-mixture likelihood bank]\label{lem:uf-mixture-bank}
The probability of a wrong bank stop is at most $2\rho\epsilon$
uniformly over $q\in\mathcal Q_\rho$, and its false-positive probability
is at most $\rho\epsilon$. These statements concern the event that the
bank stops and do not assume termination at every magnitude.

For $d\in I_j$ in a selected band,
\begin{equation}\label{eq:uf-bank-negative}
 \E_{d,-}T_{\rm bank}
 \le\frac3{4d^2}
       [1+\log(1/\epsilon)+\log(1/\pi_j)].
\end{equation}
Suppose a finite number $V(d)>0$ satisfies
\[
 \beta_\mu(n):=\int e^{-2nx^2}\,\mu(dx)\le V(d)
 \qquad\text{for every integer }n\ge d^{-2}.
\]
Then
\begin{equation}\label{eq:uf-bank-positive}
 \E_{d,+}T_{\rm bank}
 \le\frac6{d^2}
 \left[1+\log(1/\epsilon)+\log(1/\pi_j)
                   +\log_+\frac{V(d)}{\rho}\right].
\end{equation}
In particular the bank has finite conditional work on every selected
band whenever such an envelope is supplied.
\end{lemma}

\begin{proof}
A nonnegative supermartingale $Y_n$ with $Y_0=1$ has
$\Pp(\sup_nY_n\ge a)\le1/a$: stopping at the first crossing or a fixed
horizon gives $1\ge\E Y_{T\wedge m}\ge a\,\Pp(T\le m)$, and increasing
the horizon proves the claim. This elementary maximal inequality also
works when the process is a supermartingale only on supported prefixes.

Under a positive coin of gap $d$, the next multiplier of $L_{j,n}^-$
has mean $1-4r_j^-d\le1$. Its negative-crossing probability is at most
$\epsilon\pi_j$. A countable union bound gives conditional
negative-crossing probability at most $\epsilon$, so its positive-sign
contribution under any admissible prior is at most $\rho\epsilon$.

Under the full all-negative mixture, $M_n$ is a nonnegative
supermartingale on supported prefixes. If a prefix $x$ has positive
null probability, conditional expectation sums only children $xu$
with positive null probability; their numerator probabilities sum to
at most the numerator at $x$. This proves the supermartingale
inequality even when $\mu$ has mass at $d=1/2$. A zero-denominator prefix
has null probability zero. Thus a positive crossing has null-mixture
probability at most $\rho\epsilon$. Its actual negative-sign
contribution is no larger, since it multiplies the conditional
crossing probabilities by $1-q(d)\le1$. The two error claims follow.

For negative work, fix a selected-band gap $d\in I_j$ and
$r=r_j^-\in[d/2,d]$, with $r\le1/4$. The matching log likelihood has
mean increment
\[
 g(d,r)=\frac12\log(1-4r^2)
          +d\log\frac{1+2r}{1-2r}
 \ge-\frac83r^2+4dr\ge\frac43d^2.
\]
Here $\log(1-x)\ge-x/(1-x)$ and
$\log((1+x)/(1-x))\ge2x$ establish the first lower bound; the
concave quadratic takes the value $4d^2/3$ at both endpoints
$r=d/2,d$. Every upward log increment is at most $\log(3/2)<1$.
For the first correct-reference crossing $\tau$, the stopped-sum
identity at $\tau\wedge m$ gives
\[
 g(d,r)\E(\tau\wedge m)
 =\E\log L^-_{j,\tau\wedge m}
 \le1+\log(1/(\epsilon\pi_j)).
\]
Monotone convergence proves integrability of $\tau$ and the stated
work bound, since $T_{\rm bank}\le\tau$. No unbounded stopping identity
is invoked before integrability is established.

For positive work, let
$E_n=\{\widehat p_n-1/2\ge d/2\}$. On this event, for $x<1/2$,
\[
 \log\frac{P_{1/2-x,n}}{P_{{\rm f},n}}
 =\frac n2\log(1-4x^2)
   -n(\widehat p_n-\tfrac12)
                 \log\frac{1+2x}{1-2x}
 \le-2nx^2.
\]
For $x=1/2$ the likelihood is zero, since $E_n$ entails at least one
head. Therefore $P_{-,\mu,n}/P_{{\rm f},n}\le\beta_\mu(n)$.
The matching positive reference $r=r_j^+\in[d/4,d/2]$ satisfies,
on the same event,
\[
 \log\frac{P_{1/2+r,n}}{P_{{\rm f},n}}
 \ge n\left[-\frac83r^2+2dr\right]\ge nd^2/3.
\]
Set
\[
 U=1+\log(1/\epsilon)+\log(1/\pi_j)
                  +\log_+(V(d)/\rho),\qquad
 n_0=\left\lceil3d^{-2}U\right\rceil.
\]
For every $n\ge n_0$, $E_n$ implies
$M_n\ge\pi_je^{nd^2/3}/V(d)\ge1/(\rho\epsilon)$.
Thus $T_{\rm bank}>n$ implies $E_n^c$, whose positive-coin probability
is at most $e^{-nd^2/2}$. Summing the survival probabilities yields
\[
 \E_{d,+}T_{\rm bank}
 \le3d^{-2}U+2+2d^{-2}\le6d^{-2}U.
\]
The last inequality uses $U\ge1$ and $d^2\le1/4$.

For a countable bank, the positive mixture is a consistent probability
law and the negative stopping event is a measurable countable union.
At any fixed time only finitely many negative references can cross:
$L^-_{j,n}\le(3/2)^n$ requires
$\pi_j\ge1/[\epsilon(3/2)^n]$ for a crossing. All preceding arguments
therefore apply to finite and countable banks alike.
\end{proof}

\subsection{Racing procedures and retaining a genuine unit}

\begin{lemma}[Observable race]\label{lem:uf-race}
Suppose two classifiers run on the odd and even subsequences of one
unit's observation tape, with independent auxiliary seeds. Stop the
physical procedure when either classifier stops and use its output.
On completed counterfactual subsequences its work satisfies
\[
 T_{\rm race}\le2\min\{T_1,T_2\}.
\]
Its error event is contained in the union of the component wrong-stop
events. If either component has a finite deterministic cap, the race
terminates with a finite deterministic cap.
\end{lemma}

\begin{proof}
By physical time $2n$, each component has received $n$ observations.
If either component would stop after $n$ of its own observations, the
race stops no later than $2n$. This proves the pathwise work bound,
including an infinite stopping time for the other component.
The selected output is wrong only if the component that produces it
makes a wrong stop. Conditional on the latent $(D,S)$, both
subsequences are iid Bernoulli tapes of the correct mean, so all
component error and work bounds apply. The policy uses only the
physical observed prefix and its seeds.
\end{proof}

\subsection{Finite multilevel allocation at every positive budget}

\begin{lemma}[Independent Poisson allocation]\label{lem:uf-poisson}
Let $Z_0,\ldots,Z_J\in\{0,1\}$ be observable single-unit classifiers,
each with finite deterministic cap. Suppose uniformly over
$q\in\mathcal Q_\rho$ that their error probabilities are at most
$e_\ell$, their false-positive probabilities are at most $f_\ell$,
and their expected work is at most $C_\ell>0$.
Put
\[
 v_0=\rho+f_0,\qquad
 v_\ell=e_\ell+e_{\ell-1}\quad(\ell\ge1),\qquad
 k_0=C_0,\qquad k_\ell=C_\ell+C_{\ell-1}\quad(\ell\ge1).
\]
Assume $v_\ell>0$ and set $A_J=\sum_{\ell=0}^J\sqrt{v_\ell k_\ell}$.
For every $B>0$ there is an observable population estimator of
uniform expected work at most $B$ and worst RMS error at most
\begin{equation}\label{eq:uf-poisson-risk}
 \sqrt{e_J^2+A_J^2/B}.
\end{equation}
\end{lemma}

\begin{proof}
A level-zero correction uses a fresh genuine unit and returns
$\Delta_0=Z_0$. For $\ell\ge1$, a correction draws one fresh latent unit,
runs $Z_\ell$ on its observations, then runs $Z_{\ell-1}$ with a fresh
seed on the next unused observations of the same unit, and returns
$\Delta_\ell=Z_\ell-Z_{\ell-1}$.

Conditional on the latent state, the unused suffix after a finite
stopping time is an iid tape with the same Bernoulli mean: for each
possible stopped prefix, independence of its subsequent coordinates
gives the claim, and summing over the stopped prefixes preserves it.
This conditional argument is applied to successive unused blocks of
the same latent unit, even though the first block has random length.
The second classifier therefore has its original marginal distribution
under the same latent state. No new magnitude or sign is drawn between
the two classifiers. Their disagreement implies that at least one
misclassifies that common sign, and hence
\[
 \E_q\Delta_\ell^2\le v_\ell,\qquad
 \E_q W_\ell\le k_\ell,\qquad
 \sum_{\ell=0}^J\E_q\Delta_\ell=\E_qZ_J.
\]
At level zero use
$\Pp_q(Z_0=1)\le\Pp_q(S=1)+f_0\le\rho+f_0$.

Independently across levels, draw Poisson replication counts with means
\[
 n_\ell=\frac B{A_J}\sqrt{\frac{v_\ell}{k_\ell}}.
\]
For each replication use an independent genuine unit, and return
\[
 \widehat\theta
 =\sum_{\ell=0}^J\frac1{n_\ell}
       \sum_{i=1}^{N_\ell}\Delta_{\ell,i},
 \qquad N_\ell\sim\operatorname{Poisson}(n_\ell).
\]
Its mean is $\E_qZ_J$, so its absolute bias is at most $e_J$.
The expected work is at most
$\sum_\ell n_\ell k_\ell=B$.

For independent marks $\Delta$ and a Poisson count of mean $n$,
conditioning on the count gives
$\Var(\sum_{i=1}^{N}\Delta_i)=n\E\Delta^2$.
Consequently the independent level variances sum to at most
\[
 \sum_{\ell=0}^J\frac{v_\ell}{n_\ell}=\frac{A_J^2}{B}.
\]
Bias squared plus variance proves \eqref{eq:uf-poisson-risk}.
All counts are almost surely finite, all component caps are finite,
and all first and subsequent observations are included in $W_\ell$.
\end{proof}

The allocation may use Poisson means below one; this is valid for
every $B>0$. The bound is uniform over the original rarity class.
Finite-level bounds can be optimized by taking an infimum over
finite constructions without asserting the existence of a single
infinite-level algorithm.


\section{One-unit null-work obstructions}\label{sec:lower-foundations}

This appendix supplies the sign-testing and annular ingredients used in
Section~\ref{sec:matching-lower}. The weighted sign-testing mechanism is
an adaptation of Chen--Li--Qiao's gap-entropy argument, including their
Bernoulli extension \cite[Lemma 4.1 and Remark D.4]{clq}. We give the full
finite-gap version needed here, followed by the scale-sensitive gate and
its actual minimum-band-mass constraint.

\subsection{Weighted sign testing}

\begin{lemma}[Weighted sign entropy]\label{lem:weighted-sign}
Let $d_1>\cdots>d_K>0$, with $d_j/d_{j+1}\ge2$ and $d_1\le1/128$.
One observable binary sign test, with the same independent seed law at all
gaps, has error at most $1/10$ under every Bernoulli law $1/2\pm d_j$.
Suppose its stopping time $\tau$ has finite expectation under both signs at
every such gap.  Put $c_j=d_j^2\E_{-d_j}\tau$.
For every probability vector $\pi$ on these gaps,
\begin{equation}\label{eq:sc-weighted-sign}
 \sum_{j=1}^K\pi_jc_j\ge
 \frac{1+\operatorname{Ent}(\pi)}{214}.
\end{equation}
No termination or correctness assumption under the fair coin is required.
\end{lemma}

\begin{proof}
First, the ordinary stopped change-of-measure inequality gives a uniform
positive floor for each $c_j$.  The terminal sign decision has minus/plus
KL divergence at least
\[
 \operatorname{kl}(9/10,1/10)=\frac45\log9>\frac85.
\]
One Bernoulli observation has divergence at most
\[
 \frac{(2d_j)^2}{(1/2+d_j)(1/2-d_j)}\le\frac{64}{3}d_j^2.
\]
The stopped chain rule, justified by finite expected endpoint work, implies
$c_j>3/40>1/16$.

Set $b_0=1/1024$ and
\[
 n_j=\left\lfloor b_0/d_j^2\right\rfloor,\qquad
 u_j=\left\lceil4\E_{-d_j}\tau\right\rceil.
\]
Since $d_j\le1/128$, we have $n_j\ge b_0/(2d_j^2)\ge8$.
The KL divergence between the complete prefixes of length $n_j$ at the
minus and plus laws is at most $1/48$.  These prefixes include the
independent seed and unused bits if the test has already stopped.
Pinsker's inequality therefore gives total variation less than $1/8$.
An early minus is a plus-law error, so its probability under the minus law
is at most $1/10+1/8$.  Markov's inequality gives
$\Pp_{-d_j}(\tau>u_j)\le1/4$.  Consequently the event
\[
 E_j=\{\text{output minus},\ n_j<\tau\le u_j\}
\]
has $\Pp_{-d_j}(E_j)>1/4$.

Let $\Pp_0$ denote the fair-coin law with the same seed and policy.
Comparing deterministic prefixes of length $u_j$ gives KL at most
$4d_j^2u_j\le16c_j+1$.  Binary data processing and
$\operatorname{kl}(p,z)\ge p\log(1/z)-\log2$ yield
\begin{equation}\label{eq:sc-window-mass}
 \Pp_0(E_j)\ge\exp(-64c_j-8).
\end{equation}
Only deterministic finite prefixes have been compared with the fair law.

We next bound a weighted overlap pathwise.  On a path with finite terminal
time $t$, participation in $E_j$ implies
\[
 t>n_j\ge\frac{b_0}{2d_j^2},\qquad
 t\le4\E_{-d_j}\tau+1.
\]
In particular $t\ge2$ and
$c_j\ge(t-1)d_j^2/4\ge td_j^2/8$.
Order the participating gaps from smallest to largest.  Their squared
values grow by at least a factor of four, and the smallest has
$td_j^2>b_0/2$.  The corresponding costs are thus bounded below by
$b_0 4^r/16=4^{r-7}$, for $r=0,1,\ldots$.  Hence
\[
 \sum_j\frac{\1_{E_j}}{1+c_j}
 \le\sum_{r\ge0}\frac1{1+4^{r-7}}
 \le7+\sum_{r\ge7}4^{7-r}=\frac{25}{3}.
\]
On a nonterminating fair path no $E_j$ occurs, so the same inequality
holds.  Integrating and using \eqref{eq:sc-window-mass}, together with
$1+c_j\le e^{c_j}$, gives
\[
 Z:=\sum_je^{-65c_j}\le(25/3)e^8,
 \qquad \log Z<11.
\]
Gibbs' inequality for $\pi$ and the normalized masses $e^{-65c_j}/Z$
implies
\[
 \operatorname{Ent}(\pi)\le65\sum_j\pi_jc_j+\log Z.
\]
Writing $\bar c=\sum_j\pi_jc_j$ and retaining the stronger floor
$\bar c\ge3/40$, we obtain
\[
 1+\operatorname{Ent}(\pi)
 \le65\bar c+9+\log(25/3)
 \le\left[65+\frac{40}{3}\{9+\log(25/3)\}\right]\bar c
 <214\bar c.
\]
The windows need not be disjoint; only the displayed weighted-overlap
bound is used.
\end{proof}


\subsection{A moment comparison on the half-prior class}

Write $X_\rho=\{0\le q\le1/2:\int q\,d\mu\le\rho\}$.
For an integrable observable one-unit output $Y$ with finite uniform
variance and uniform bias at most $\rho/64$, put $V=\sup_{q\in X_\rho}\Var_qY$ and
$M=\sup_{q\in X_\rho}\E_qY^2$.

\begin{lemma}[Moment comparison]\label{lem:lf-moment}
For $0<\rho\le1/4$, one has $V\ge2\rho/3$ and $M\le3V/2$.
\end{lemma}
\begin{proof}
Let $f(d)=\E_{d,+}Y-\E_{d,-}Y$. Comparing the null population with
the constant prior $q=\rho$ gives $\int f\,d\mu\ge31/32$.
Conditional variance at this constant prior and Jensen's inequality yield
\[
 V\ge\rho(1-\rho)\int f^2\,d\mu
 \ge\rho(1-\rho)(31/32)^2\ge2\rho/3.
\]
Every mean has absolute value at most $65\rho/64$, so
$M\le V+(65\rho/64)^2\le3V/2$.
The constant prior also gives integrability of both fixed-sign second
moments; all conditional statements may be taken outside a common null set.
\end{proof}

\subsection{An activation gate that respects the magnitude scale}

The new lower-bound ingredient charges activation at the information
scale \(a^2\E_0T\).  In this subsection \(\mu\) is any probability law
supported on \((0,a]\), \(a\le1/4\), and
\[
 \mathcal X_\rho=
 \{q:0\le q\le\tfrac12,\ \textstyle\int q\,d\mu\le\rho\}.
\]
Let \((Y,T)\) be an observable one-coin procedure with uniform bias at
most \(\rho/64\), finite uniform variance \(V\), and finite expected work
under every \(q\in\mathcal X_\rho\).  Write
\[
 M_0=\sup_{q\in\mathcal X_\rho}\E_qY^2,\qquad
 W_0=\E_0T,\qquad P=M_0W_0.
\]
Here $M_0$ is the uniform second-moment bound, whereas $W_0$ is
null work; the subscript on $M_0$ does not restrict its supremum to the null.
The moment comparison in Lemma~\ref{lem:lf-moment} gives
\(M_0\le3V/2\).

\begin{lemma}[A scale-sensitive baseline and gated activation]
\label{lem:information-gate}
Put
\[
 C_g=e^7,\qquad \kappa=64/3,\qquad
 c_g=(256\kappa C_g^4)^{-1}.
\]
Every procedure above satisfies
\begin{equation}
 P\ge\frac{\rho}{64a^2}.
 \label{eq:cross-baseline}
\end{equation}
If \(a^2P\le c_g\), there is an observable procedure \(Z\) on the same
one-coin experiment, with moment bound \(m_Z\), null work \(W_0^Z\), and
a measurable set \(\mathcal G\) of good magnitudes, such that
\begin{equation}
 \begin{gathered}
 \sup_{q\in\mathcal X_\rho}\E_qZ^2\le m_Z,\qquad
 \rho/16\le m_Z\le2C_g^4a^2P,\\
 m_ZW_0^Z\le C_g^4P,\qquad
 \mu(\mathcal G^c)\le144m_Z.
 \end{gathered}
 \label{eq:cross-gate-bounds}
\end{equation}
At every good gap, both signs have finite conditional expected work and
\begin{equation}
 \left|\E[Z\mid d,0]\right|\le\frac18,\qquad
 \E[Z\mid d,1]\ge\frac12,\qquad
 \E[Z^2\mid d,0]\le\frac{m_Z}{\rho},\qquad
 \E[Z^2\mid d,1]\le\frac{3m_Z}{\rho}.
 \label{eq:cross-gate-good}
\end{equation}
One common sign test with both errors below \(1/10\) on these gaps uses
\begin{equation}
 n=\left\lceil1024m_Z/\rho\right\rceil
 \le1040m_Z/\rho
 \label{eq:cross-gate-restarts}
\end{equation}
independent restarts of \(Z\) on unused observations of the same coin.
Restart independence is conditional on its fixed latent magnitude and
sign, not after averaging over the population.
Its negative expected work is \(n\,\E[T_Z\mid d,0]\).
\end{lemma}

\begin{proof}
We first prove \eqref{eq:cross-baseline}.  Let \(P_-\) and \(P_+\) be
the full transcript laws with uniformly negative and uniformly positive
signs, respectively.  The priors \(q=0\) and \(q\equiv\rho\) give
\[
 \E_+Y-\E_-Y\ge31/32,\qquad
 \E_-Y^2\le M_0,\qquad \E_+Y^2\le M_0/\rho.
\]
The positive moment and expected work are finite because \(q\equiv\rho\)
is admissible.  Revealing \(D\) for the information calculation, and
using the stopped chain rule, gives
\[
 h(P_-,P_+)\le\tfrac12\kappa a^2W_0,
\]
where \(h=1-\operatorname{Aff}\).
Indeed the one-observation endpoint divergence is at most
\(\kappa d^2\) for \(d\le1/4\).
For any square-integrable output \(F\),
\begin{equation}
 |\E_PF-\E_QF|^2
 \le4(\E_PF^2+\E_QF^2)h(P,Q).
 \label{eq:cross-hellinger-moment}
\end{equation}
Thus
\[
 (31/32)^2
 \le(5\kappa/2)a^2P/\rho,
\]
which implies \eqref{eq:cross-baseline}.

The decision to take a first observation is seed-only and precedes
all data. The information gate below is instead an observable
finite-prefix stopping rule; neither step reveals a latent magnitude.
For the gate, on a full observation tape define
\[
 S_n=\sum_{i=1}^n(2X_i-1),\qquad
 n_0=\lceil a^{-2}\rceil,\qquad H_0=\lceil a^{-1}\rceil,
\]
\begin{equation}
 G=\min\{n_0,\inf\{n\ge1:|S_n|\ge H_0\}\},\qquad
 \sigma=T\wedge G,\qquad E=\{T\ge G\}.
 \label{eq:cross-gate}
\end{equation}
The event \(E\) is measurable at \(\sigma\), including the tie \(T=G\).
Include the entire independent seed in the stopped prefix.
At \(\sigma\), \(|S_\sigma|\le H_0\) and \(\sigma\le n_0\).
For a fixed \(d\le a\) and sign \(s\), the likelihood relative to the
fair stopped-prefix law is
\[
 \exp\left\{(2s-1)S_\sigma\operatorname{arctanh}(2d)
       +\frac{\sigma}{2}\log(1-4d^2)\right\}.
\]
The inequalities
\[
 \operatorname{arctanh}(2d)\le8a/3,\qquad
 -\tfrac12\log(1-4d^2)\le8a^2/3
\]
show that the absolute logarithm is at most \(37/6<7\).
Every fixed-state prefix law, and the all-negative magnitude mixture,
therefore has density in \([C_g^{-1},C_g]\) relative to fair.

Let \(p=\Pp_{\rm fair}(E)\) and \(p_s(d)=\Pp_{d,s}(E)\).  Then
\begin{equation}
 p/C_g\le p_s(d)\le C_gp.
 \label{eq:cross-gate-probabilities}
\end{equation}
Bounded fair stopping gives
\(\E_{\rm fair}S_\sigma^2=\E_{\rm fair}\sigma\).
A gate activation either reaches \(n_0\) or hits \(H_0\), so
\begin{equation}
 p\le\left(n_0^{-1}+H_0^{-2}\right)\E_{\rm fair}\sigma
 \le2a^2\E_{\rm fair}\sigma
 \le2C_ga^2W_0.
 \label{eq:cross-gate-cost}
\end{equation}

Set \(Y_e=Y\1_{\{T<G\}}\), explicitly zero on \(E\).
This is a function of the bounded prefix.
Its second moment under either fixed state is at most \(C_g^2M_0\).
The stopped endpoint divergence and the same prefix domination give
\[
 \KL(P_{-d}^{\sigma},P_{+d}^{\sigma})
 \le\kappa d^2\E_{-d}\sigma
 \le\kappa C_g^2a^2W_0.
\]
Since \(h\le\KL/2\), \eqref{eq:cross-hellinger-moment} implies
\begin{equation}
 |\E_{+d}Y_e-\E_{-d}Y_e|^2
 \le4\kappa C_g^4a^2P.
 \label{eq:cross-early-output}
\end{equation}
Under \(a^2P\le c_g\), every such early mean difference is at most
\(1/8\) in absolute value.  In this branch \(p>0\): if \(p=0\), the
gate never activates under any state, so \(Y=Y_e\), contradicting
the integrated mean separation \(31/32\) from \(q=0,q\equiv\rho\).

Here is an actual observable procedure realizing the conditioning.
Keep the same genuine coin.  Repeatedly run the original procedure,
using a fresh original seed and unused iid flips for each trial.
Reject a trial if its original stopping time is before its gate.
When a trial reaches the gate, accept it, continue that original trial
to its original stopping time, and return
\[
 Z=C_gpY.
\]
Every rejected observation is charged.  Conditional on \(d,s\), trials
are iid with acceptance probability \(p_s(d)>0\).
Writing \(w_s(d)=\E[T\mid d,s]\), the renewal identity gives
\begin{equation}
 \E[T_Z\mid d,s]=\frac{w_s(d)}{p_s(d)}.
 \label{eq:cross-renewal}
\end{equation}
This identity includes the work of all rejected trials and the final
accepted trial.  Moreover,
\[
 \E[Z^2\mid d,s]
 =\frac{C_g^2p^2}{p_s(d)}\E[Y^2\1_E\mid d,s]
 \le C_g^3p\,\E[Y^2\mid d,s].
\]
With \(m_Z=C_g^3pM_0\), integration and
\eqref{eq:cross-gate-cost}--\eqref{eq:cross-renewal} give
\begin{equation}
 \sup_q\E_qZ^2\le m_Z,\qquad
 m_Z\le2C_g^4a^2P,\qquad
 W_0^Z\le C_gW_0/p,\qquad
 m_ZW_0^Z\le C_g^4P.
 \label{eq:cross-gate-moments}
\end{equation}
Expected work under every \(q\) is finite by the same comparison.

We establish conditional sign separation directly; conditioning has
not preserved the estimator's means.  Write
\[
 f_s(d)=\E[Y\mid d,s],\quad
 \ell_s(d)=\E[Y\1_E\mid d,s],\quad
 z_s(d)=\E[Z\mid d,s].
\]
The set \(\{f_+-f_-<3/4\}\) has mass below \(2\rho\).
Otherwise a fractional prior of total mass \(\rho\) on that set would
contradict the bias guarantee and the mean separation \(31\rho/32\).
Outside that set, \eqref{eq:cross-early-output} gives
\(\ell_+-\ell_-\ge5/8\).
Also
\[
 z_s(d)=\frac{C_gp}{p_s(d)}\ell_s(d),\qquad
 1\le\frac{C_gp}{p_s(d)}\le C_g^2.
\]
If \(|z_-(d)|\le1/8\), then \(|\ell_-(d)|\le1/8\),
\(\ell_+(d)\ge1/2\), and \(z_+(d)\ge1/2\).
The exceptional mass for \(|z_-|>1/8\) is at most \(64m_Z\).
Since
\[
 m_Z\le2C_g^4c_g=1/(128\kappa)<1/256,
\]
this preliminary mean-good set has mass at least
\(1-2\rho-64m_Z\ge1/4\).
Taking \(q\equiv\rho\) in the second-moment bound yields
\(m_Z\ge\rho/16\).

Finally, put \(a_Z(d)=\E[Z^2\mid d,0]\) and
\(b_Z(d)=\E[Z^2\mid d,1]\).
The set \(\{a_Z>m_Z/\rho\}\) has mass at most \(\rho\).
The set \(\{b_Z-a_Z>2m_Z/\rho\}\) has mass below \(2\rho\),
again by using a fractional admissible prior on a set of mass
at least \(2\rho\).  Removing these sets and null sets of infinite
conditional work gives total exceptional mass
\[
 5\rho+64m_Z\le144m_Z
\]
and proves \eqref{eq:cross-gate-good}.
Average \(n\) independent restarts of \(Z\) on unused blocks of the
same coin, and threshold at \(5/16\).
Both mean margins are at least \(3/16\), so Chebyshev gives errors
at most
\[
 \frac{3m_Z/\rho}{n(3/16)^2}\le1/12<1/10.
\]
The bound on \(n\) follows from \(m_Z\ge\rho/16\), and work adds,
proving \eqref{eq:cross-gate-restarts}.
\end{proof}

\subsection{A scaled annular obstruction}

For a positive vector \(v=(v_i)\), set
\[
 \mathcal G(v)=A[1+\Ent(v/A)],\qquad A=\sum_i v_i.
\]
A useful elementary identity is
\begin{equation}
 \frac{\partial\mathcal G}{\partial v_i}
 =1+\log(A/v_i)\ge1.
 \label{eq:cross-entropy-monotone}
\end{equation}
Thus this cost-entropy functional is coordinatewise increasing.
In particular, if \(u\le v_i\le U\) for \(K\) coordinates, then
\begin{equation}
 uK(1+\log K)\le\mathcal G(v)\le UK(1+\log K).
 \label{eq:cross-entropy-sandwich}
\end{equation}

\begin{proposition}[A gated annular lower bound]
\label{prop:scaled-annular}
Suppose \(\mu\) is supported on \((0,a]\), \(a\le1/4\).
Take a nonempty finite annular family
\[
 I_j=(b_j/2,b_j],\qquad b_j/b_{j+1}\ge4,\qquad
 b_1\le1/128,
\]
with positive masses \(w_j=\mu(I_j)\).  Define
\[
 w_*=\min_jw_j,\qquad
 A=\sum_jw_j/b_j^2,\qquad
 \pi_j=\frac{w_j/b_j^2}{A},\qquad
 \Gamma^{\circ}=A[1+\Ent(\pi)].
\]
Every one-coin procedure with the assumptions above satisfies
\begin{equation}
 V\E_0T\ge
 \frac{e^{-28}}{2\,000\,000}
       \min\{\rho\Gamma^{\circ},w_*/a^2\}.
 \label{eq:cross-scaled-annular}
\end{equation}
\end{proposition}

\begin{proof}
If \(a^2P\ge c_g\), use \(w_*\le1\) and \(M_0\le3V/2\).
Otherwise apply Lemma~\ref{lem:information-gate}.
If \(m_Z>w_*/288\), then
\[
 P>w_*/(576C_g^4a^2).
\]
If \(m_Z\le w_*/288\), at least half the mass of every annulus is
good.  Averaging selects \(d_j\in I_j\cap\mathcal G\) with
\[
 w_-^Z(d_j)\le\frac4{w_j}\int_{I_j}w_-^Z(d)\,\mu(dd).
\]
The selected gaps are separated by a factor greater than two and
bounded by \(1/128\).
Apply Lemma~\ref{lem:weighted-sign} to the common restart test,
with weights \(\pi_j\):
\[
 \frac{1+\Ent(\pi)}{214}
 \le n\sum_j\pi_jd_j^2w_-^Z(d_j)
 \le\frac{4n}{A}W_0^Z.
\]
Consequently,
\[
 \rho\Gamma^{\circ}
 \le856\rho nW_0^Z
 \le890\,240\,m_ZW_0^Z
 \le890\,240\,C_g^4P.
\]
Converting \(P=M_0W_0\) to \(VW_0\) completes all branches with the
displayed conservative constant.  No atomlessness is used.
\end{proof}

The entropy branch is available when $\rho a^2\Gamma^{\circ}\le w_*$.
This condition is invariant under a common rescaling of all magnitudes.

\section{Convex comparison, common pairs, and additive localization}
\label{sec:transfer-foundations}

This appendix proves the specialized transfer inputs used by the scalar
lower bound and the stopping-program comparison. The finite-experiment
convex-estimation principle has antecedents in
Juditsky--Nemirovski and Polyanskiy--Wu \cite{jn,pw};
the midpoint adaptive decomposition is due to Lee--Valiant \cite{lv}.
The arguments needed for this model are included in full.

\subsection{Finite categorical convex comparison}

\begin{lemma}[Finite categorical convex comparison]
\label{lem:finite-comparison}
Let $X$ be a nonempty compact convex subset of a finite-dimensional real vector
space. Let $\ell:X\to\R$ be affine, and let $p(x)$ be a strictly positive
categorical probability vector with affine continuous coordinates. For $n\ge1$
independent observations from $p(x)$, define
\[
 \Omega_n=\max\{\ell(x)-\ell(y):x,y\in X,
                  \operatorname{Aff}(p(x),p(y))^n\ge1/2\}.
\]
For every $\eta>0$, there is a bounded additive estimator
$F=c+\sum_{i=1}^n\phi_i(Z_i)$ satisfying
\begin{equation}
 \sup_{x\in X}\{\E_x(F-\ell(x))^2\}^{1/2}
 \le\tfrac32\Omega_n+\eta.
 \label{eq:finite-comparison}
\end{equation}
\end{lemma}

\begin{proof}
We reproduce the short argument needed for the numerical constant. Set
$s=\log2$, let $\mathcal F$ be the finite-dimensional space of additive
functions, and define, for $\alpha>0$,
\[
 \Phi(x,y;\phi,\alpha)=\ell(x)-\ell(y)
 +\alpha\log\E_y e^{\phi/\alpha}
 +\alpha\log\E_x e^{-\phi/\alpha}+2\alpha s.
\]
The function is continuous and concave in $(x,y)$: independence writes each
log moment as a sum of logarithms of positive affine forms. It is convex in
$(\phi,\alpha)$ by the perspective transformation of a log moment generating
function. Sion's minimax theorem \cite{sion} therefore applies to $X\times X$ and
$\mathcal F\times(0,\infty)$. Write
$2S=\inf_{\phi,\alpha>0}\max_{x,y}\Phi$.

For a fixed pair, the half log likelihood ratio belongs to $\mathcal F$ and
Cauchy--Schwarz gives
\[
 \inf_{\psi\in\mathcal F}
 \bigl[\log\E_y e^\psi+\log\E_x e^{-\psi}\bigr]
 =2\log\operatorname{Aff}(p(x)^{\otimes n},p(y)^{\otimes n}).
\]
Substituting $\phi=\alpha\psi$ and minimizing over $\alpha>0$ gives the value
$\ell(x)-\ell(y)$ when the affinity is at least $1/2$, and $-\infty$ otherwise.
Consequently $2S=\Omega_n$. The constrained maximum exists because its feasible
set is nonempty and compact.

For $R>S$, choose $(\phi,\alpha)$ with $\max\Phi\le2R$. On the diagonal,
Cauchy--Schwarz implies $\Phi(x,x;\phi,\alpha)\ge2\alpha s$, hence
$\alpha\le R/s$. Set
\[
 U=\max_x\{\ell(x)+\alpha\log\E_xe^{-\phi/\alpha}\},\qquad
 V=\max_x\{-\ell(x)+\alpha\log\E_xe^{\phi/\alpha}\}.
\]
Then $U+V+2\alpha s\le2R$. The estimator
$F=\phi+(U-V)/2$ satisfies
\[
 \E_xe^{\pm(F-\ell(x))/\alpha}\le e^{R/\alpha-s}.
\]
Therefore $\Pp_x\{|F-\ell(x)|>R+t\}\le e^{-t/\alpha}$ for $t\ge0$.
Tail integration yields
\[
 \E_x(F-\ell(x))^2
 \le R^2+2R\alpha+2\alpha^2
 \le R^2\left(1+\frac2{\log2}+\frac2{(\log2)^2}\right)<9R^2.
\]
Let $R\downarrow S$, with the slack chosen for $\eta$. Every component is
bounded because its alphabet is finite. This proves
\eqref{eq:finite-comparison}.
\end{proof}


\subsection{Midpoint decomposition for interleaved observations}

We use $h(P,Q)=1-\Aff(P,Q)$, so $0\le h\le1$ and
$\operatorname{TV}(P,Q)\le\sqrt{2h(P,Q)}$.
The following finite-prefix argument is the midpoint decomposition of
Lee--Valiant \cite[Lemmas 17--18]{lv}, written in the form needed here.

\begin{lemma}[Midpoint work decomposition]\label{lem:tf-midpoint}
Fix two endpoint priors $q_-,q_+$ and their midpoint
$\bar q=(q_-+q_+)/2$. Suppose every capped randomized one-unit policy
$A$ obeys
\[
 h(P_-^A,P_+^A)\le\kappa\,\E_{\bar q}T_A.
\]
Then every population algorithm with a deterministic total cap satisfies
\[
 h(P_-^{\rm pop},P_+^{\rm pop})
 \le\kappa\,\E_{\bar q}T_{\rm pop}.
\]
It suffices that the local hypothesis hold up to that same total cap.
If instead $q_\pm\le1/2$ and the local bound uses null work
$\lambda\E_0T_A$, the conclusion holds with $\kappa=2\lambda$.
The statement also holds for finitely many observable component types,
each with its own endpoint pair and midpoint, when the same coefficient
bounds every local policy for every type. Types may be selected adaptively, but each fresh unit's type is fixed
before its first observation and retained on every revisit.
\end{lemma}

\begin{proof}
Fix total cap $K$ and condition on the algorithm's independent seed.
Defer latent draws until a unit is first observed, which changes no
observable law. Index units by that first observation. A consistent global leaf $z$
contains local binary words $h_i(z)$, with empty words for unused slots.
The endpoint leaf densities factor as
\[
 p_\pm^{\rm pop}(z)=\prod_{i=1}^K p_\pm(h_i(z)).
\]
The common policy-consistency indicator is implicit. Put
$m_i=(p_-(h_i)+p_+(h_i))/2=p_{\bar q}(h_i)$ and
$c_i=\sqrt{p_-(h_i)p_+(h_i)}$. Then
\[
 \prod_i m_i-\prod_i c_i
 \le\sum_i(m_i-c_i)\prod_{j\ne i}m_j.
\]
The integral of $\prod_i m_i$ over consistent leaves is one, since it
is the population law in which every fresh unit has midpoint prior.
Integrating the inequality therefore bounds population squared Hellinger
by a sum of hybrid squared Hellinger distances. In hybrid $i$ only
unit $i$ has an endpoint prior; all other units have midpoint priors.

Simulate all other units from independent midpoint tapes included in
an auxiliary seed, and run the population rule with one genuine unit
in slot $i$. This is a capped local policy. Its full transcript,
including that seed, determines the hybrid global transcript.
Data processing bounds hybrid distance by local distance. The expected
midpoint local works sum to $\E_{\bar q}T_{\rm pop}$ because each
purchased observation belongs to exactly one slot. This proves the
first assertion, including after integrating the original seed.
The charged expectation is under the midpoint population, admissible
by convexity. It is not replaced by the average endpoint work, which
can be much smaller under an adaptive stopping rule.

If $q_\pm\le1/2$, the midpoint latent law dominates half the null
latent law. The same measure inequality holds after passing through
any observation policy and auxiliary seed. Thus
$\E_0T_A\le2\E_{\bar q}T_A$.

For observable types, use the type selected at a slot in each leaf.
The same product factorization holds, with its type-specific local
kernel. To realize a hybrid locally, first simulate the midpoint
history before the slot's first observation; its type and that history are
part of the independent seed. The genuine unit has the selected type's
prescribed latent law independently of that seed. Future simulated midpoint
units use fresh independent seeds. Conditional on the preceding history, the hypothesis
for the selected type applies. Sum and integrate as above. No latent
magnitude or sign is supplied to the local policy.
\end{proof}

For an unbounded stopping rule, apply the conclusion to every bounded
global prefix. Total variation of the full transcript is the increasing
limit of the prefix total variations: bounded cylinder functions generate
the transcript sigma-field, or equivalently one applies conditional
expectation to the likelihood densities under the dominating sum law.
Finite expected endpoint work ensures almost-sure finite termination.
A common pair valid for every cap therefore gives the same terminal
testing bound. Equivalently, couple the stopped experiment with its
cap-$L$ truncation under each endpoint. Their total coupling error is at
most $\Pp_-(T>L)+\Pp_+(T>L)$, which tends to zero. This passage
requires no common deterministic bound on stopping.

\subsection{The half-prior class and null work}

Put
\[
 X_\rho=\left\{q:0\le q\le\tfrac12,\quad
                  \theta(q):=\int q\,d\mu\le\rho\right\},
 \qquad 0<\rho\le\tfrac14.
\]
For a capped observable one-coin procedure with bounded real output \(Y\),
write
\[
 V=\sup_{q\in X_\rho}\Var_qY,\qquad
 M=\sup_{q\in X_\rho}\E_qY^2,\qquad W_0=\E_0T,
\]
and define
\begin{equation}
 \mathfrak J_0(b)=\inf\left\{VW_0:
       \sup_{q\in X_\rho}|\E_qY-\theta(q)|\le b,
       \ T\text{ capped},\ Y\text{ bounded}\right\}.
 \label{eq:null-bias-complexity}
\end{equation}
An empty infimum is infinity.  The function is nonincreasing in \(b\),
and \(\mathfrak J_0(\rho/2)=0\).

\subsection{Transfer at an arbitrary bias tolerance}

\begin{theorem}[Bias-indexed common pairs and population lower bounds]
\label{thm:bias-transfer}
Suppose that every capped bounded one-coin output with uniform bias at
most \(b>0\) over \(X_\rho\) satisfies \(VW_0\ge J\), where
\(0<J<\infty\).  Then, for every \(0<\delta\le b/3\), there is one pair
\(q_-,q_+\in X_\rho\) with target separation at least \(\delta\) such
that simultaneously for all capped randomized local policies \(A\),
\begin{equation}
 h(P_{q_-}^A,P_{q_+}^A)
 \le42\frac{\delta^2}{J}\E_0T_A.
 \label{eq:bias-common-pair}
\end{equation}
Consequently,
\begin{equation}
 R_{\mu,\rho}(B)\ge\frac1{128}
                 \min\{b,\sqrt{J/B}\}.
 \label{eq:bias-population-transfer}
\end{equation}
The conclusion includes fully adaptive creation, revisiting, interleaving,
nonlinear outputs and unbounded stopping.  
\end{theorem}

\begin{proof}
The premise forces \(b<\rho/2\), since the zero-work constant \(\rho/2\)
would otherwise contradict it.  First fix a finite strictly positive
categorical statistic of a capped policy of null work \(w>0\).
For \(n\) independent runs, let
\(F=c+\sum_{i=1}^nY_i\) be a bounded additive estimator of worst RMS
\(r\).  Put \(b_i=\E_0Y_i\), \(b_F=\E_0F\), and choose an independent
uniform index \(I\).  Running just its component and returning
\(Z=n(Y_I-b_I)+b_F\) preserves all population means at null work \(w\).

The laws under \(q\) and zero share a subprobability of mass
\(1-\theta(q)\), obtained by weighting the negative-sign kernel by
\(1-q(d)\).  The elementary common-mixture variance inequality gives
\[
 (\E_qY_i-b_i)^2\le
 \frac{2\theta(q)}{1-\theta(q)}
                 (\Var_qY_i+\Var_0Y_i).
\]
Indeed, either mean's displacement from the normalized common part is
at most the square root of \(\theta/(1-\theta)\) times its standard
deviation; square the sum of the two displacements.  Independence in
\(F\) now yields
\[
 \sup_q\Var_qZ\le
 nr^2\left(1+\frac{4\rho}{1-\rho}\right)\le\tfrac73nr^2.
\]
If \(r<b\), apply the assumed obstruction to \(Z\); if \(r\ge b\),
retain that inequality.  In both cases,
\[
 r\ge\min\{b,\sqrt{3J/(7nw)}\}.
\]

Lemma~\ref{lem:finite-comparison} gives an additive estimator of
RMS at most \((3/2)\Omega_n+\eta\), where \(\Omega_n\) is the maximum
target separation whose product affinity is at least \(1/2\).
Letting \(\eta\downarrow0\) therefore gives
\[
 \Omega_n\ge\tfrac23\min\{b,\sqrt{3J/(7nw)}\}.
\]
Set \(x=J/(21\delta^2w)\).  If \(x\ge2\), take
\(n=\lfloor x\rfloor\); then \(\Omega_n\ge2\delta\), and the
corresponding pair satisfies
\[
 h\le\frac{\log2}{n}
 \le42\log2\,\frac{\delta^2w}{J}
 \le42\frac{\delta^2w}{J}.
\]
If \(x<2\), the last right side exceeds one, so any separated pair
suffices.  If \(w=0\), the seed-only first-observation event has
probability zero; all local laws agree and \(q_-=0,q_+=\delta\)
suffices.

We now obtain one pair simultaneously for every capped local policy.
For an arbitrary capped policy, partition its full transcript into finitely
many measurable cells, then replace the cell independently by a uniform
cell with probability \(\varepsilon>0\).  Run the policy even on
replacement outcomes, so its null work is unchanged.  The resulting
strictly positive category probabilities and the target have compact
convex finite-dimensional range, so the preceding comparison applies.
For any finite collection of these experiments, independently mix them
and reveal the chosen index.  Both work and squared Hellinger distance
are affine in the mixture weights.  Hellinger is continuous and convex
in the parameter pair.  Finite-simplex minimax gives a pair satisfying
every constraint in the collection.  The constraint sets are closed in
the weak-star compact convex set
\[
 \{(q_-,q_+)\in X_\rho^2:
                 \theta(q_+)-\theta(q_-)\ge\delta\}.
\]
The finite intersection property supplies one pair for all experiments.
For fixed $b$ and $\delta$, this pair is chosen before the adaptive
population algorithm and its total-observation cap. The same pair
works for every capped local policy and every cap.
Let \(\varepsilon\downarrow0\), then use that full squared Hellinger
distance is the supremum over finite measurable partitions.  This proves
\eqref{eq:bias-common-pair} for arbitrary transcript spaces.

The midpoint decomposition in Lemma~\ref{lem:tf-midpoint}
bounds population squared Hellinger by
\(84\delta^2B/J\): the local null cost is at most twice the midpoint
cost, and summing midpoint costs counts each observation once.  Set
\[
 \delta=\frac1{32}\min\{b,\sqrt{J/B}\}.
\]
The distance is at most \(84/1024<1/8\), hence TV is at most \(1/2\).
The same pair works for every deterministic observation cap.  Finite
expected endpoint work and increasing capped prefixes give the same
terminal TV bound for unbounded stopping.  Thresholding an estimator at
the target midpoint gives RMS at least \(\delta/4\), proving
\eqref{eq:bias-population-transfer}.  A budget over the full rarity class
implies the required budget over \(X_\rho\).
\end{proof}

\subsection{A direct variational null-work obstruction}

For Bernoulli observations, define the known variational modulus \cite{repo}
\begin{equation}
 \Psi_{\mu,\rho}(h)=
 \sup_{\substack{0\le w\le1\\\int w\,d\mu\le\rho}}
 \int\min\{w(d),h\sqrt{w(d)}/d\}\,\mu(dd).
 \label{eq:inherited-psi}
\end{equation}
The proof below establishes the exact null-work input required here
directly from fixed-sign comparisons. No optimizer formula is needed.

\begin{proposition}[Variational accuracy-dependent null obstruction]
\label{prop:psi-null}
For every probability measure \(\mu\) on \((0,1/2]\) and every
\(0\le b\le\rho/64\),
\begin{equation}
 \mathfrak J_0(b)\ge
 \sup_{h>0}\frac{(\Psi_{\mu,\rho}(h)-4b)_+^2}{192h^2}.
 \label{eq:psi-null-product}
\end{equation}
No nonatomicity, finite support, or separation of magnitudes is assumed.
\end{proposition}

\begin{proof}
Fix an eligible \((Y,T)\), and define its conditional response, moments,
and negative work by
\[
 f(d)=\E_{d,+}Y-\E_{d,-}Y,\qquad
 A_s(d)=\E_{d,s}Y^2,\qquad w_-(d)=\E_{d,-}T.
\]
We first show
\begin{equation}
 f(d)^2\le64d^2w_-(d)[A_-(d)+A_+(d)].
 \label{eq:conditional-response-cost}
\end{equation}
For \(d\le1/4\), the stopped KL chain rule gives
\(\KL(P_{d,-}^A,P_{d,+}^A)\le(64/3)d^2w_-(d)\).
For any square-integrable output under two laws,
\[
 |\E_PY-\E_QY|^2
 \le4(\E_PY^2+\E_QY^2)h(P,Q),\qquad
 h(P,Q)\le\tfrac12\KL(P,Q).
\]
These prove \eqref{eq:conditional-response-cost} with coefficient
\(128/3<64\).  For \(d>1/4\), the conditional laws agree whenever
the seed requests no observation.  Their Hellinger distance is therefore
at most the activation probability, hence at most \(w_-(d)\).
The same mean inequality gives
\(f(d)^2\le4w_-(d)(A_-+A_+)\le64d^2w_-(d)(A_-+A_+)\).
This also handles \(d=1/2\) without a singular endpoint KL comparison.

For every feasible \(w\) in \eqref{eq:inherited-psi}, the allocation
\(q=w/2\) belongs to \(X_\rho\), and
\[
 \int w(A_-+A_+)\,d\mu
 \le2\int[(1-w/2)A_-+(w/2)A_+]\,d\mu\le2M.
\]
Weighted Cauchy--Schwarz and \eqref{eq:conditional-response-cost} yield
\begin{equation}
 \left(\int\frac{\sqrt w\,|f|}{d}\,d\mu\right)^2
 \le W_0\int\frac{wf^2}{d^2w_-}\,d\mu
 \le128MW_0.
 \label{eq:weighted-response-bound}
\end{equation}
On \(\{w_-=0\}\), the response inequality forces \(f=0\), and the
quotient is zero.  Truncation justifies Cauchy--Schwarz before passing to
the limit if an integral is initially extended-valued.

Put \(v_h=\min\{w,h\sqrt w/d\}\).  Comparing the mean under the
admissible allocation \(q=v_h/2\) with the null mean, and using both
bias bounds, gives
\[
 \int v_h f\,d\mu\ge\int v_h\,d\mu-4b.
\]
Since \(v_h\le h\sqrt w/d\),
\eqref{eq:weighted-response-bound} implies
\[
 MW_0\ge\frac{(\int v_h\,d\mu-4b)_+^2}{128h^2}.
\]
Take the supremum over \(w\).  To convert the second moment to variance,
the null and constant-prior populations force the all-positive and
all-negative means to differ by at least \(31/32\).  The mixture with
constant sign prior \(\rho\) consequently has variance at least
\(\rho(1-\rho)(31/32)^2\ge2\rho/3\).  Also every admissible mean has
absolute value at most \(65\rho/64\).  Thus
\(M\le V+(65\rho/64)^2\le3V/2\), proving
\eqref{eq:psi-null-product}.
\end{proof}

This product bound accounts for accuracy-dependent tail difficulty in the
same framework as the annular scale-entropy obstruction.  At a fixed gap
\(d\), the directly evaluated identity
\(\Psi(h)=\min\{\rho,h\sqrt\rho/d\}\) recovers product scale
\(\rho/d^2\) as \(b\downarrow0\).  For the quadratic law it supplies
the inverse-square logarithm; the additional scale-entropy factor still
requires the annular argument.



\begin{proposition}[An optimizer for the variational modulus]
\label{prop:psi-optimizer}
For $0<\rho<1$, there is a unique $\alpha>0$ satisfying
\[
 \int\min\{1,\alpha^2/d^2\}\,\mu(dd)=\rho.
\]
The allocation $w_*(d)=\min\{1,\alpha^2/d^2\}$ maximizes
$\Psi_{\mu,\rho}(h)$ simultaneously for every $h>0$. In particular,
$\Psi(h)=\rho$ when $h\ge\alpha$. Fractional allocation at an atom is
permitted and may be essential; this is distinct from the hard
assignment statement for the entropy-selection infimum.
\end{proposition}
\begin{proof}
The mass on the left is continuous in $\alpha$, tends to zero at zero
and to one at infinity, and is strictly increasing while it is below
one. These assertions follow by dominated convergence and the positive
mass of unsaturated magnitudes. Thus the stated root exists uniquely.
If $h\ge\alpha$, the allocation attains the upper bound $\rho$.
If $h<\alpha$, put $\lambda=h/(2\alpha)<1/2$. The function
\[
 \phi_d(w)=\min\{w,(h/d)\sqrt w\},\qquad 0\le w\le1,
\]
is concave. For $d>\alpha$, the point $w_*=(\alpha/d)^2$ is in its
square-root branch and has derivative $\lambda$. For $h<d\le\alpha$,
the left derivative at $w_*=1$ is $h/(2d)\ge\lambda$; for $d\le h$,
$\phi_d(w)=w$. Hence $w_*$ maximizes $\phi_d(w)-\lambda w$
pointwise. Integrating and using its mass $\rho$ proves optimality over
all allocations of mass at most $\rho$, including atomic laws.
\end{proof}

\begin{remark}[The bias subtraction cannot be uniformly reduced]
Let $\rho=1/8$, $\varepsilon=2^{-10}$, and
$\mu=(1-\varepsilon)\delta_{1/2}+\varepsilon\delta_d$.
The one-flip output $Y=X-b$, with $b=\varepsilon(1-2d)/4$, has bias
at most $b$ on $X_\rho$, null work one, and variance at most $1/4$.
Indeed its bias is
$\varepsilon(1/2-d)(1-2q(d))-b\in[-b,b]$.
Yet $\Psi(d)\ge\varepsilon$. Replacing $4b$ by $cb$ with any fixed
$c<4$ in \eqref{eq:psi-null-product} would make its right-hand side
diverge as $d\downarrow0$. The coefficient four cannot be reduced in
this form of the bound, even if its denominator constant is changed.
\end{remark}

\subsection{A direct sum of conditional experiments}

Let $\sigma_1,\ldots,\sigma_m$ be nonzero finite submeasures with
$\sum_i\sigma_i\le\mu$. Write $\tau_i=\sigma_i((0,1/2])$ and
$\nu_i=\sigma_i/\tau_i$. Disjoint restrictions are a special case;
fractional restrictions of atoms are also allowed. Choose
\[
 0<r_i\le1/4,\qquad \sum_i\tau_i r_i\le\rho.
\]
Within component $i$, use the half-prior subclass
\[
 X_{r_i}(\nu_i)=\{0\le q\le1/2:\ \textstyle\int q\,d\nu_i\le r_i\}.
\]
The original population lower bound is valid on this subclass of the full
rarity class.

\begin{theorem}[Additive localization]\label{thm:additive-localization}
Suppose that, for each $i$, every capped bounded observable one-unit output
with uniform bias at most $b_i>0$ over $X_{r_i}(\nu_i)$ satisfies
\[
 \left(\sup_q\Var_qY\right)\E_{i,0}T\ge J_i>0.
\]
Then, for every $B>0$,
\begin{equation}\label{eq:additive-localization}
 \boxed{\quad
 R_{\mu,\rho}(B)\ge\frac1{128}
 \sum_{i=1}^m\tau_i\min\{b_i,\sqrt{J_i/B}\}.
 \quad}
\end{equation}
The same total budget $B$ occurs in every summand. The component label
is chosen before the first observation and retained on every revisit;
its conditional magnitude and sign remain hidden.
\end{theorem}

\begin{proof}
Theorem~\ref{thm:bias-transfer} applies within every component. Set
\[
 \delta_i=\frac1{32}\min\{b_i,\sqrt{J_i/B}\}.
\]
It supplies a pair $q_{i,-},q_{i,+}\in X_{r_i}(\nu_i)$, oriented to have
conditional target separation at least $\delta_i$, such that simultaneously
for every capped local policy $A$,
\begin{equation}\label{eq:additive-local-pair}
 h(P_{i,-}^A,P_{i,+}^A)
 \le42\frac{\delta_i^2}{J_i}\E_{i,0}T_A
 \le\frac{42}{1024B}\E_{i,0}T_A.
\end{equation}

Consider a stronger experiment that can select a component at no cost and
then receive a genuine coin with magnitude law $\nu_i$. Its component label
is visible, while its magnitude and sign remain hidden. The algorithm can
adaptively select, revisit, and interleave these coins. Under the conditional
midpoint prior, $q\le1/2$, so the midpoint law dominates half the null law.
Lemma~\ref{lem:tf-midpoint} therefore bounds terminal squared
Hellinger by $84/1024<1/8$: the local coefficient in
\eqref{eq:additive-local-pair} is common to all components, and the summed
conditional midpoint costs count each purchased observation once.
Increasing capped prefixes handles unbounded stopping. The total target
separation is at least $\sum_i\tau_i\delta_i$.

Every original algorithm can be simulated in this stronger experiment.
Privately choose the component label with probabilities $\tau_i$ whenever
it creates a fresh coin. On the remaining probability mass, simulate the
known all-negative residual law for free. Keep the chosen label when the
coin is revisited. For fractional submeasures, integrating this private
mixture label gives the original admissible parameters
\[
 q_\pm(d)=\sum_i\frac{d\sigma_i}{d\mu}(d)q_{i,\pm}(d),
 \qquad 0\le q_\pm\le1/2,\qquad
 \int q_\pm\,d\mu\le\sum_i\tau_i r_i\le\rho.
\]
The number of genuine oracle observations does not exceed the original
number. Thus the oracle lower bound applies to the original experiment.
The Hellinger bound gives TV at most $1/2$, and midpoint testing gives RMS
at least one quarter of the target separation. This is
\eqref{eq:additive-localization}.
\end{proof}

For a direct explicit input, suppose component $i$ is supported in
$(0,a_i]$, $a_i\le1/4$, and select a finite separated family of small-gap
annuli. Let $\Gamma_i^{\circ}$ and $w_{*,i}$ denote its conditional cost entropy
and minimum band mass. Proposition~\ref{prop:scaled-annular} gives
\[
 b_i=r_i/64,\qquad
 J_i=c_{\mathrm{ann}}
       \min\{r_i\Gamma_i^{\circ},w_{*,i}/a_i^2\},\qquad
 c_{\mathrm{ann}}=\frac{e^{-28}}{2{,}000{,}000}.
\]
Write the corresponding unnormalized quantities with respect to
$\sigma_i$ as $\Gamma_{\sigma_i}^{\circ}=\tau_i\Gamma_i^{\circ}$ and
$w_{*,\sigma_i}=\tau_iw_{*,i}$. The conversion uses the exact identity, with $\theta=\tau r$,
\[
 \tau^2\min\{r\Gamma^{\circ},w_*/a^2\}
 =\min\{\theta\Gamma_{\sigma}^{\circ},\tau w_{*,\sigma}/a^2\}.
\]
In particular the second branch retains the factor $\tau$. Consequently,
\begin{equation}\label{eq:additive-annular}
 R_{\mu,\rho}(B)\gtrsim
 \sum_i\min\left\{\tau_i r_i,
 \sqrt{\frac{
 \min\{\tau_i r_i\Gamma_{\sigma_i}^{\circ},
           \tau_iw_{*,\sigma_i}/a_i^2\}}B}\right\}.
\end{equation}
Supremizing this expression over finite decompositions and rarity
allocations gives an explicit lower functional. The bias-indexed form
\eqref{eq:additive-localization} also permits each component's accuracy
scale to be chosen separately. Section~\ref{sec:matching-lower} applies this expression to components
selected by the entropy--tail optimizer.


\subsection{A universal risk floor}

\begin{lemma}[Bernoulli-label floor]\label{lem:tf-floor}
For every magnitude law, $0<\rho\le1/4$, and $B>0$,
\[
 R_{\mu,\rho}(B)\ge\frac1{16}
       \min\{\rho,\sqrt{\rho/B}\}.
\]
\end{lemma}
\begin{proof}
Restrict to constant sign allocations. Give an oracle the exact sign
of a unit when its first observation is requested, for one charged
observation. It can sample $D\sim\mu$
independently and simulate all future flips for free, so every original
algorithm can be simulated with no greater work. The informative labels
are iid Bernoulli with the unknown constant prior.

Take $p_0=\rho/2$ and $p_1=p_0+\delta$, where
$\delta=\min\{\rho/2,\sqrt{\rho/B}/4\}$. Both are admissible and
$p_1(1-p_1)\ge3\rho/8$. The stopped KL chain rule, first at bounded
prefixes and then by monotone convergence, gives
\[
 \KL(P_0,P_1)\le B\,\operatorname{kl}(p_0,p_1)
 \le\frac{8B\delta^2}{3\rho}\le\frac16.
\]
Thus $\operatorname{TV}(P_0,P_1)\le1/2$. Thresholding any estimator at
the target midpoint implies worst RMS at least $\delta/4$.
Finally
$\delta/4=\min\{\rho/8,\sqrt{\rho/B}/16\}
\ge\min\{\rho,\sqrt{\rho/B}\}/16$.
\end{proof}

\section{Representation and realization of the stopping flow}
\label{sec:flow-foundations}

This section justifies the normalized flow cone and robust program in
\eqref{eq:dual-flow-cone}--\eqref{eq:dual-primal}. Triangular encodings
of stopped coin experiments have direct antecedents in Lee--Valiant
\cite{lv}; occupation-measure formulations of optimal stopping are
also classical \cite{cs}. The proofs below give the representation
and estimator realization needed for this particular robust program.

\subsection{The support function of the rarity class}

For bounded measurable $f:(0,1/2]\to\R$, put
\[
 Q_{\mu,\rho}(f)
 =\sup_{q\in\mathcal Q_\rho}\int q(d)f(d)\,\mu(dd).
\]

\begin{lemma}[A fractional upper-tail formula]\label{lem:ff-support}
The support function satisfies the attained identity
\begin{equation}\label{eq:ff-support}
 Q_{\mu,\rho}(f)
 =\min_{\ell\ge0}
       \left\{\rho\ell+\int(f-\ell)_+\,d\mu\right\}.
\end{equation}
The minimum can be restricted to
$[0,\|f_+\|_\infty]$. The supremum on the left is attained by selecting
the largest positive values of $f$, allowing a fractional selection
on the cutoff level set.
\end{lemma}

\begin{proof}
For every $0\le q\le1$ and $\ell\ge0$,
$qf\le q\ell+(f-\ell)_+$. Integration and $\int q\,d\mu\le\rho$
prove the upper bound by the right side.

If $\mu(f>0)\le\rho$, take $q=\1_{\{f>0\}}$ and $\ell=0$.
Otherwise choose a positive cutoff $\ell$ satisfying
\[
 \mu(f>\ell)\le\rho\le\mu(f\ge\ell).
\]
Such a cutoff exists by the monotonicity and one-sided continuity of
the tail distribution. Take $q=1$ on $\{f>\ell\}$, $q=0$ on
$\{f<\ell\}$, and choose a constant fraction in $[0,1]$ on
$\{f=\ell\}$ so that its total mass is $\rho$.
The resulting value is
\[
 \int_{\{f>\ell\}}f\,d\mu
   +\ell[\rho-\mu(f>\ell)]
 =\rho\ell+\int(f-\ell)_+\,d\mu.
\]
This proves equality and attainment, including atoms of $\mu$ or
atoms in the distribution of $f$. Cutoffs larger than
$\|f_+\|_\infty$ cannot improve on the endpoint.
\end{proof}

For an array $u=(u_{nk})$, define
\[
 L_u^\pm(d)
 =\sum_{n,k}u_{nk}\binom nk
          (1/2\pm d)^k(1/2\mp d)^{n-k},
 \qquad
 I_u^-=\int L_u^-\,d\mu,\qquad
 F_u=L_u^+-L_u^-.
\]
The root or final-depth terms are included only when the corresponding
array permits them. The convention $0^0=1$ covers the deterministic
Bernoulli endpoints. By the definition of $m(q)$,
\begin{equation}\label{eq:ff-robust-linear}
 \sup_{q\in\mathcal Q_\rho}\langle m(q),u\rangle
 =I_u^-+Q_{\mu,\rho}(F_u).
\end{equation}
In particular, the two bias suprema in \eqref{eq:dual-primal} are
\[
 c_0+I_z^-+Q_{\mu,\rho}(F_z-1),
 \qquad
 -c_0-I_z^-+Q_{\mu,\rho}(1-F_z).
\]
The moment and work suprema use \eqref{eq:ff-robust-linear} with
$u=t$ and $u=c$. Thus the robust program requires only integration
against the known magnitude law and a scalar cutoff optimization.

\subsection{Arbitrary seeds and ordered histories aggregate to the cone}

\begin{proposition}[Normalized prefix aggregation]
\label{prop:ff-aggregation}
Every capped one-unit policy with $1\le T\le N$, arbitrary independent
auxiliary randomization, and bounded terminal output has normalized
arrays in the cone \eqref{eq:dual-flow-cone} with root intensity one.
Its mean, raw second moment, and expected work under any $q$ are
respectively
\[
 \langle m(q),z\rangle,\qquad
 \langle m(q),t\rangle,\qquad
 \langle m(q),c\rangle.
\]
Multiplying the four arrays by $\lambda\ge0$ gives the corresponding
intensity-$\lambda$ representation. The stopping representation alone
requires no output or bounded-output assumption.
\end{proposition}

\begin{proof}
Place all auxiliary randomization in one seed $r$, drawn from a law
$\nu$ independent of the latent unit and its observations. Once the
seed is fixed, all decisions and outputs are deterministic functions
of the ordered observed prefix.

For each ordered binary prefix $x$ of length $n$, let $A_x$ be the set
of seeds for which the policy reaches $x$ and stops there; let $C_x$
be the set for which it reaches $x$ and continues. These sets are
disjoint. For $n\ge1$ their union is exactly $C_{x^-}$, where $x^-$
is the prefix with its last symbol removed. Write
\[
 \alpha_x=\nu(A_x),\quad
 \gamma_x=\nu(C_x),\quad
 \zeta_x=\int_{A_x}Y(r,x)\,\nu(dr),\quad
 \tau_x=\int_{A_x}Y(r,x)^2\,\nu(dr).
\]
At the root, $\gamma_\varnothing=1$ and $\alpha_\varnothing=0$;
at depth $N$, $\gamma_x=0$.

For $s=(n,k)$, sum these quantities over all prefixes with $k$ heads
and divide by $\binom nk$:
\[
 a_{nk}=\frac{\sum_{|x|=n,\,\#1(x)=k}\alpha_x}{\binom nk},
 \qquad
 c_{nk}=\frac{\sum_{|x|=n,\,\#1(x)=k}\gamma_x}{\binom nk},
\]
and analogously define $z_{nk}$ and $t_{nk}$ from $\zeta_x,\tau_x$.
Aggregating $\alpha_x+\gamma_x=\gamma_{x^-}$ gives
\[
 a_{nk}+c_{nk}
 =\frac{n-k}{n}c_{n-1,k}
    +\frac{k}{n}c_{n-1,k-1}.
\]
Cauchy--Schwarz on the disjoint union of the prefix-indexed
stopping-seed sets
gives $z_{nk}^2\le a_{nk}t_{nk}$. Hence the arrays belong to the
stated cone.

Conditional on Bernoulli mean $p$, every ordered prefix in this
count state has the same probability $p^k(1-p)^{n-k}$.
Summing its output integrals therefore gives the stated moment
identities after mixing over the latent state. Expected work is the
sum of the probabilities of continuing at all prefixes, including
the root: each continuation requests exactly one charged observation.
This proves the work identity and accounts for the first observation.

Multiplication by $\lambda$ preserves the linear flow identities and
the homogeneous perspective inequalities. The argument also applies
to outputs with finite first and second seed aggregates. For the
bounded outputs used in finite constructions, these are automatic.
\end{proof}

The normalized arrays encode probabilities over the independent seed
for fixed observation prefixes. They are not the unconditional
probabilities of reaching a count state. This distinction is the
reason the Pascal coefficients in \eqref{eq:dual-flow-cone} are
$(n-k)/n$ and $k/n$.

\subsection{Every feasible point is an observable estimator}

\begin{proposition}[Compound-Poisson realization]
\label{prop:ff-realization}
Every feasible point of \eqref{eq:dual-primal} constructs an observable
population estimator with uniform expected work at most $B$ and
worst RMS error at most $\sqrt{b^2+v}$. In particular,
\[
 R_{\mu,\rho}(B)\le P_N(B).
\]
\end{proposition}

\begin{proof}
Suppose first that $\lambda>0$. Draw a Poisson number of independent
genuine units with mean $\lambda$. Each unit first requests one
observation. At an observed state $(n,k)$, stop with probability
$a_{nk}/(a_{nk}+c_{nk})$ and continue otherwise. At a stopping state
with $a_{nk}>0$, return $z_{nk}/a_{nk}$. Unreachable states may have
arbitrary actions, and outputs at zero stopping flow can be set to
zero. The constraint $c_{Nk}=0$ enforces the cap.

For a fixed Bernoulli mean $p$, the expected number of units reaching
$(n,k)$ is
\[
 (a_{nk}+c_{nk})\binom nk p^k(1-p)^{n-k}.
\]
This follows by induction from the normalized Pascal recursion.
For example, a continuation at $(n-1,k)$ followed by a zero
contributes
$c_{n-1,k}\binom{n-1}{k}p^k(1-p)^{n-k}$, and the other parent
contributes the analogous term with $k-1$. Their sum is exactly the
displayed expression. Multiplying by the stopping or continuation
probability gives the same formula with $a_{nk}$ or $c_{nk}$ alone.

Return $c_0$ plus the sum of all terminal unit outputs. Independence
of the genuine units and the compound-Poisson moment identities give
\begin{align}
 \E_q\widehat\theta
 &=c_0+\langle m(q),z\rangle,\label{eq:ff-realized-mean}\\
 \Var_q(\widehat\theta)
 &=\left\langle m(q),\frac{z^2}{a}\right\rangle
 \le\langle m(q),t\rangle,\label{eq:ff-realized-variance}\\
 \E_qT&=\langle m(q),c\rangle.\label{eq:ff-realized-work}
\end{align}
Set $z_{nk}^2/a_{nk}=0$ where $a_{nk}=0$; the perspective constraint
then also gives $z_{nk}=0$. Every sum is finite, and each nonzero
terminal output is a finite real number. The root contribution to
\eqref{eq:ff-realized-work} is $\lambda$, since $m_{00}(q)=1$.
The work constraint therefore charges every first observation and
implies $\lambda\le B$.

The bias, moment, and work inequalities in the program prove the
claimed uniform guarantees. When $\lambda=0$, the flow recursion
forces $a=c=0$ and the perspective constraint forces $z=0$.
The deterministic estimator $c_0$ realizes the required bound;
any unused nonnegative moment slack is harmless.
\end{proof}

Allowing zero-observation marks would not improve the attainable
Poisson guarantee. Indeed, the event that a one-unit procedure stops
before its first observation depends only on its independent seed.
Poisson thinning separates those marks from the active units. Replace
their total random contribution by its deterministic expectation.
This preserves every population mean and the work, while removing
their nonnegative variance contribution. The remaining active
intensity, together with the adjusted offset $c_0$, has precisely
the root convention in \eqref{eq:dual-flow-cone}.

\subsection{Finite-support laws give finite second-order cone programs}

Suppose $\mu=\sum_{i=1}^{m}w_i\delta_{d_i}$ is supplied by its positive
masses and distinct positive magnitudes. For each of the four support
functions in the two bias, moment, and work constraints, introduce
one cutoff $\ell\ge0$ and nonnegative epigraph variables $r_i$ with
\[
 r_i\ge f(d_i)-\ell.
\]
The corresponding robust term is represented by
$\rho\ell+\sum_iw_ir_i$, by Lemma~\ref{lem:ff-support}.
All $f(d_i)$ are affine functions of the relevant flow or output
array. Thus these epigraph constraints are linear.

The perspective constraints are second-order cones:
\[
 z_{nk}^2\le a_{nk}t_{nk}
 \quad\Longleftrightarrow\quad
 \|(2z_{nk},a_{nk}-t_{nk})\|_2\le a_{nk}+t_{nk},
 \qquad a_{nk},t_{nk}\ge0.
\]
One additional cone represents an epigraph variable $u\ge b^2$,
for example
$\|(2b,u-1)\|_2\le u+1$, and the objective is $u+v$.

Consequently the finite-support version has $O(N^2+m)$ scalar
variables and cone or linear constraints, with a dense coefficient
description of size $O(mN^2)$ after the Bernoulli kernels are
evaluated. This is a dimension count at a specified cap, without a
uniform claim about the sufficient cap, numerical conditioning, or
bit complexity.


\section{An exact dual for the robust stopping program}
\label{sec:exact-dual}

Triangular stopping representations and their adaptive comparison have
close antecedents in Lee--Valiant \cite{lv}. Here the sign allocation
ranges over the full convex rarity class, and bias, variance, and work
are robust constraints. We derive the exact four-prior dual of this
program and absorb its variance and work references into the endpoints.
Appendix~\ref{sec:flow-foundations} proves the stopping-flow representation
and its observable realization in full.

Throughout this section, $\mu$ is any probability measure on $(0,1/2]$,
$0<\rho\le1/4$, $B>0$, and $N\ge1$ is an integer. Put
\[
 \mathcal Q_\rho=\{q\in L^\infty(\mu):0\le q\le1,\
                \theta(q)=\textstyle\int q\,d\mu\le\rho\}.
\]
For a state $s=(n,k)$, $0\le k\le n\le N$, define
\begin{align}
 m_{nk}(q)&=\int\binom nk\bigl[
 (1-q(d))(1/2-d)^k(1/2+d)^{n-k}\notag\\[-2pt]
 &\hspace{39mm}+q(d)(1/2+d)^k(1/2-d)^{n-k}\bigr]\,\mu(dd).
 \label{eq:dual-count-coefficients}
\end{align}
In particular $m_{00}(q)=1$. These are probabilities of count states at
fixed depth, before imposing a stopping rule.

\subsection{The primal and its four-prior dual}

For clarity, the finite program being dualized is recalled in compact
form. Its nonnegative stopping and continuation arrays $a,c$, real output
array $z$, and nonnegative moment array $t$ obey
\begin{align}
 &a_{00}=z_{00}=t_{00}=0,\qquad c_{00}=\lambda\ge0,
       \qquad c_{Nk}=0,\notag\\
 &a_{nk}+c_{nk}=\frac{n-k}{n}c_{n-1,k}
                    +\frac{k}{n}c_{n-1,k-1},
       \qquad 1\le n\le N,\notag\\
 &z_{nk}^2\le a_{nk}t_{nk}.
 \label{eq:dual-flow-cone}
\end{align}
Missing indices are zero. These are normalized prefix-weight flows, not
unconditional state probabilities. Write $\mathcal C_N$ for this closed
convex cone. The squared program value is
\begin{equation}
\begin{split}
 P_N(B)^2=\min\;&b^2+v\quad
       \text{over }(a,c,z,t,\lambda)\in\mathcal C_N,\ c_0\in\R,\ b,v\ge0,\\
 \text{subject to }&
 \sup_{q\in\mathcal Q_\rho}
       \bigl|c_0+\langle m(q),z\rangle-\theta(q)\bigr|\le b,\\
 &\sup_q\langle m(q),t\rangle\le v,\qquad
   \sup_q\langle m(q),c\rangle\le B.
\end{split}
\label{eq:dual-primal}
\end{equation}
Stopping sums exclude the root and continuation sums exclude depth $N$.
Every feasible point realizes a Poisson estimator with uniform RMS at
most $\sqrt{b^2+v}$ and expected work at most $B$
(Proposition~\ref{prop:ff-realization}).

Use the closed perspective
\[
 \mathfrak r(x,y)=
 \begin{cases}
 x^2/y,&y>0,\\
 0,&x=y=0,\\
 +\infty,&y=0, x\ne0.
 \end{cases}
\]
Given four priors $q_L,q_H,q_v,q_w\in\mathcal Q_\rho$, put
$d_s=m_s(q_H)-m_s(q_L)$, $v_s=m_s(q_v)$, $w_s=m_s(q_w)$, and define
\begin{equation}
 K_N(q_L,q_H;q_v,q_w)
 =\sup_{A:\,1\le T_A\le N}
 \frac{\displaystyle\sum_s a_s^A\mathfrak r(d_s,v_s)}
      {\displaystyle\sum_s c_s^A w_s}.
 \label{eq:dual-information-ratio}
\end{equation}
Here $A$ is a one-unit policy that always requests its first observation;
its normalized flow has root intensity one. A zero stopping coefficient
contributes zero even when the perspective is infinite. The denominator
is $\E_{q_w}T_A\in[1,N]$.

For a signed measure $P_H-P_L$, write
\[
 \mathcal I(P_H,P_L;P_v)
 =\int\left(\frac{d(P_H-P_L)}{dP_v}\right)^2\,dP_v,
\]
with value infinity unless the signed measure is absolutely continuous.
The numerator in \eqref{eq:dual-information-ratio} equals
$\mathcal I(P_H^A,P_L^A;P_v^A)$. Include the independent seed in the
transcript: its likelihood ratio depends only on the terminal count
state, so this identity covers order-dependent stopping and randomized
policies as well.

\begin{theorem}[Exact four-prior dual]\label{thm:exact-dual}
The minimum in \eqref{eq:dual-primal} is attained, and
\begin{equation}
 \boxed{
 P_N(B)^2=
 \max_{q_L,q_H,q_v,q_w\in\mathcal Q_\rho}
 \frac{[\theta(q_H)-\theta(q_L)]_+^2}
      {4+B K_N(q_L,q_H;q_v,q_w)}.}
 \label{eq:exact-four-prior-dual}
\end{equation}
The ratio is zero when $K_N=\infty$. The maximum is attained with positive
value, positive target separation, and finite $K_N$. Swapping the endpoints
allows the positive part to be replaced by an absolute value.
\end{theorem}

\begin{proof}
\emph{Compactness and constraint qualification.}
The finite moment image
$\{(\theta(q),m(q)):q\in\mathcal Q_\rho\}$ is compact and convex:
$\mathcal Q_\rho$ is weak-star compact, and every coordinate is a bounded
linear functional of $q$. All robust functions in the program are thus
finite continuous convex functions.

The zero-flow output $\rho/2$ gives objective $\rho^2/4$. On this objective
sublevel, $b\le\rho/2$, $v\le\rho^2/4$, and the root work charge gives
$\lambda\le B$. The normalized recursion implies $0\le a_s,c_s\le B$.
For the constant prior $q_*=\rho/2$, $m_s(q_*)>0$ unless state $s$ is
impossible under every prior. At every nonnull stopping state,
\[
 t_s\le\frac{\rho^2}{4m_s(q_*)},\qquad |z_s|\le\sqrt{Bt_s}.
\]
At universally null states set $z_s=t_s=0$, without changing any mean,
moment, or cost. The bias constraint at $q=0$ then bounds $c_0$.
Consequently this feasible sublevel, with null output coordinates removed,
is compact. The primal minimum is attained.

There is also a strict relative-interior feasible point. Choose
$0<\lambda<B/4$, continue with probability $1/2$ at each nonterminal
observed state, and stop at depth $N$. All free stopping and continuation
flows are positive and expected work is less than $2\lambda<B$.
Set $z=0$, take all free $t_s>0$, and choose sufficiently large $b,v$.
Every perspective inequality and all four robust inequalities are strict.
Finite-dimensional convex Slater duality therefore gives no duality gap,
attained inequality multipliers, and KKT conditions at a primal optimum.

\emph{Reduction to four priors.}
Split the absolute bias constraint into its two one-sided inequalities.
Each robust function is a maximum of affine functions over a compact
convex set. A subgradient is a convex combination of active affine
gradients. Since $q\mapsto(\theta(q),m(q))$ is affine and
$\mathcal Q_\rho$ is convex, this combination is represented by one active
prior. Thus the two bias subgradients, the variance subgradient, and the
work subgradient may be represented respectively by
$q_L,q_H,q_v,q_w$.

Let their nonnegative multipliers be $\alpha_L,\alpha_H,\beta,\eta$.
The free offset forces $\alpha_L=\alpha_H=\alpha$. Minimizing
$b^2-2\alpha b$ over $b\ge0$ gives $-\alpha^2$, and minimizing
$(1-\beta)v$ over $v\ge0$ requires $\beta\le1$. The remaining expression
on the flow cone is
\begin{equation}
 \beta\langle m(q_v),t\rangle
 -\alpha\langle m(q_H)-m(q_L),z\rangle
 +\eta\langle m(q_w),c\rangle.
 \label{eq:dual-cone-functional}
\end{equation}
Homogeneity requires this expression to be nonnegative everywhere on
$\mathcal C_N$; otherwise scaling a negative point gives infimum minus
infinity. When it is nonnegative its infimum is zero. The dual objective is
\begin{equation}
 \alpha[\theta(q_H)-\theta(q_L)]-\alpha^2-\eta B.
 \label{eq:dual-price-objective}
\end{equation}

We may take $\beta=1$. Indeed, for any $q_0\in\mathcal Q_\rho$, replace
$\beta m(q_v)$ by
\[
 m\bigl(\beta q_v+(1-\beta)q_0\bigr)
 =\beta m(q_v)+(1-\beta)m(q_0).
\]
The additional coefficient on $t$ is nonnegative, so cone feasibility is
preserved while the dual objective is unchanged. This also preserves
attainment.

\emph{Eliminating outputs and the work price.}
For $\alpha>0$, the perspective constraint gives, at fixed stopping flow,
\[
 \inf_{z_s^2\le a_st_s}
       \{v_st_s-\alpha d_sz_s\}
 =-\frac{\alpha^2a_s}{4}\mathfrak r(d_s,v_s).
\]
For positive $v_s$ this is completion of the square. If $a_s=0$ the
contribution is zero; if $a_s>0$, $v_s=0$, and $d_s\ne0$, the infimum is
minus infinity. The remaining zero case has value zero.
Every nonzero stopping/continuation flow scales to root intensity one.
At $\lambda=0$, $a=c=z=0$ and surplus $t$ has a nonnegative contribution.
Thus nonnegativity of
\eqref{eq:dual-cone-functional} is exactly
\[
 \eta\ge\alpha^2 K_N/4.
\]
For finite $K_N$, optimizing \eqref{eq:dual-price-objective} yields
\[
 \sup_{\alpha\ge0}
 \left\{\alpha\Delta\theta
    -\alpha^2(1+BK_N/4)\right\}
 =\frac{(\Delta\theta)_+^2}{4+BK_N}.
\]
When $K_N=\infty$, only the zero-value choice $\alpha=\eta=0$ remains.
Weak duality applies to every quadruple, and the attained KKT dual above
supplies an optimal quadruple. This proves the identity.

Finally take $q_L=0$ and $q_H=q_v=\rho$ constant. The divergence of the
latent sign contrast relative to $q_v$ is $\rho/(1-\rho)$. Data processing
bounds every stopped-transcript numerator by this quantity, and work is
at least one. Thus, for any $q_w$,
\begin{equation}
 P_N(B)^2\ge\frac{\rho^2}{4+B\rho/(1-\rho)}>0.
 \label{eq:dual-positive-value}
\end{equation}
An attained optimal dual therefore has $\alpha>0$, finite $K_N$, and
positive target separation, as asserted.
\end{proof}

\subsection{A Bellman recursion with explicit null-state conventions}

Divide each coefficient in \eqref{eq:dual-count-coefficients} by
$\binom nk$ and write $p_q(n,k)$ for the probability of one particular
ordered prefix with that count. These coefficients satisfy
\[
 p_q(n,k)=p_q(n+1,k)+p_q(n+1,k+1).
\]
In this subsection let $d_s=p_H(s)-p_L(s)$, $v_s=p_v(s)$,
$w_s=p_w(s)$, and $r_s=\mathfrak r(d_s,v_s)$ denote the corresponding
\emph{prefix} coefficients. At the root $v_0=w_0=1$ and $d_0=r_0=0$.

\begin{proposition}[Bellman representation]\label{prop:dual-bellman}
Suppose $K_N<\infty$. For $\kappa\ge0$, set
\begin{align}
 U_{Nk}(\kappa)&=r_{Nk},\notag\\
 U_{nk}(\kappa)&=\max\{r_{nk},
 U_{n+1,k}(\kappa)+U_{n+1,k+1}(\kappa)-\kappa w_{nk}\},
 \quad 1\le n<N.
 \label{eq:dual-bellman}
\end{align}
Then
\begin{equation}
 \Phi_N(\kappa):=U_{10}(\kappa)+U_{11}(\kappa)-\kappa
 =\max_{1\le T_A\le N}
 \{\mathcal I(P_H^A,P_L^A;P_v^A)-\kappa\E_{q_w}T_A\}.
 \label{eq:dual-price-value}
\end{equation}
It is finite, continuous, piecewise affine, and strictly decreasing, and
its unique nonnegative zero is $K_N$. A deterministic count-state policy
attains the maximum at every price. If some prefix has $v_s=0$ and
$d_s\ne0$, then $K_N=\infty$ instead.
\end{proposition}

\begin{proof}
At any observed prefix, stopping earns $r_s$ and continuation earns the
two child values minus the work charge $\kappa w_s$. Backward induction
therefore proves \eqref{eq:dual-bellman}. Root activation is compulsory,
which gives the final charge $\kappa$. All coefficients depend only on
the count, so the maximizing choice may be deterministic and depend only
on $(n,k)$. Randomized or order-dependent choices cannot improve it.
There are finitely many such deterministic policies, each with work in
$[1,N]$. Their affine price values have slopes in $[-N,-1]$, proving the
continuity and strict monotonicity claims. The zero-price value is
nonnegative, the value tends to minus infinity, and the zero is precisely
the maximal information-to-work ratio. If a prefix has zero reference
likelihood and nonzero contrast, the policy stopping at its depth has
infinite information.
\end{proof}

The recursion also admits an information-increment form. Put
\begin{equation}
 j_{nk}=r_{n+1,k}+r_{n+1,k+1}-r_{nk}\ge0.
 \label{eq:dual-information-increment}
\end{equation}
Nonnegativity is Cauchy--Schwarz. At a zero-reference state, finite $K_N$
forces both its contrast and its children's contrasts to vanish, so
$j_s=0$ there. If $b_s^A$ is the unnormalized continuation-prefix flow,
telescoping the flow identities gives
\begin{equation}
 \mathcal I(P_H^A,P_L^A;P_v^A)=\sum_{n<N,k}b_s^A j_s,
 \qquad \E_{q_w}T_A=\sum_{n<N,k}b_s^A w_s.
 \label{eq:dual-increment-telescope}
\end{equation}
Subtracting $r_s$ from the Bellman values gives terminal values $V_N=0$
and
\[
 V_s(\kappa)=\max\{0,j_s-\kappa w_s+V_{s0}(\kappa)+V_{s1}(\kappa)\},
 \qquad
 \Phi_N(\kappa)=j_0-\kappa+V_{10}(\kappa)+V_{11}(\kappa).
\]
Thus the ratio accounts for both information gained at a state and the
work needed to reach it. The pointwise inequalities $j_s\le\kappa w_s$
are sufficient, though generally not necessary, for $K_N\le\kappa$.

When $q_v$ dominates $q_w$ on all relevant prefixes, the posterior form is
\[
 \Phi_N(\kappa)=\sup_{1\le T\le N}
 \E_{q_v}\left[L_T^2-\kappa\sum_{n<T}W_n\right],
 \qquad L_n=d_n/v_n,\quad W_n=w_n/v_n.
\]
The process $L_n$ is a signed-likelihood martingale, initially zero. The
unnormalized Bellman recursion itself requires no such domination
assumption, and is the appropriate form at degenerate states.

\section{A midpoint modulus and a sharper adaptive comparison}
\label{sec:midpoint-dual}

Define the two-prior quantity
\begin{equation}
 D_N(B)^2=
 \sup_{q_-,q_+\in\mathcal Q_\rho}
 \frac{|\theta(q_+)-\theta(q_-)|^2}
 {4+B K_N(q_-,q_+;\bar q,\bar q)},
 \qquad \bar q=(q_-+q_+)/2.
 \label{eq:dual-midpoint-modulus}
\end{equation}

\begin{theorem}[Absorbing the two references]\label{thm:midpoint-modulus}
For every law, budget, and finite cap in the stated range,
\begin{equation}
 D_N(B)\le P_N(B)\le3D_N(B).
 \label{eq:dual-midpoint-comparison}
\end{equation}
\end{theorem}

\begin{proof}
The first inequality restricts the four-prior dual. For the second, start
with any quadruple of finite $K_N$ and set
\[
 q'_-=(q_L+q_v+q_w)/3,\qquad q'_+=(q_H+q_v+q_w)/3.
\]
Their midpoint dominates $q_v$ and $q_w$ in the following precise
experiment sense: for every policy,
$P_{\bar q'}^A\ge P_{q_v}^A/3$ and
$\E_{\bar q'}T_A\ge\E_{q_w}T_A/3$.
The signed transcript contrast is the old contrast divided by three.
Consequently its information numerator is at most one third of the old
numerator, and
\[
 K_N(q'_-,q'_+;\bar q',\bar q')
 \le K_N(q_L,q_H;q_v,q_w).
\]
The target gap is divided by three. Substitution in
\eqref{eq:dual-midpoint-modulus} and then
\eqref{eq:exact-four-prior-dual} proves the result.
\end{proof}

For midpoint references, the numerator is twice triangular
discrimination. With $h(P,Q)=1-\operatorname{Aff}(P,Q)$,
\begin{equation}
 4h(P_+,P_-)
 \le 2\int\frac{(p_+-p_-)^2}{p_++p_-}
 =\mathcal I(P_+,P_-;(P_++P_-)/2)
 \le8h(P_+,P_-).
 \label{eq:dual-hellinger-comparison}
\end{equation}
These inequalities follow by factoring
$(p_+-p_-)^2=(\sqrt{p_+}-\sqrt{p_-})^2
(\sqrt{p_+}+\sqrt{p_-})^2$ and bounding the second factor divided by
$p_++p_-$ between one and two. Hence the midpoint modulus is also a
Hellinger-information-per-work modulus, up to universal constants. It
still involves an optimal local stopping rule; it is not a closed
tail-moment or entropy formula.

\begin{lemma}[Squared-loss testing by affinity]\label{lem:affinity-testing}
If two terminal laws have target gap $\Delta$, every estimator has worst
RMS at least $\Delta\operatorname{Aff}(P_-,P_+)/2$.
\end{lemma}

\begin{proof}
The maximum MSE is at least the average MSE for equally likely endpoints.
The minimum average MSE is
\[
 \frac{\Delta^2}{2}\int\frac{p_-p_+}{p_-+p_+}.
\]
Cauchy--Schwarz gives
$\operatorname{Aff}(P_-,P_+)^2
\le2\int p_-p_+/(p_-+p_+)$, proving the claim.
\end{proof}

Let $R_N^{\mathrm{tot}}(B)$ denote population minimax RMS with at most
$N$ total observations and uniform expected work at most $B$. This total
cap is different from the per-unit cap used in $P_N$.

\begin{theorem}[A direct adaptive comparison]\label{thm:dual-adaptive}
\begin{equation}
 R_N^{\mathrm{tot}}(B)\ge\frac38D_N(B),\qquad
 P_N(B)\le8R_N^{\mathrm{tot}}(B).
 \label{eq:dual-capped-adaptive}
\end{equation}
\end{theorem}

\begin{proof}
Fix a pair in \eqref{eq:dual-midpoint-modulus}, with target gap $\Delta$
and information ratio $K$. Interpolate its endpoints toward their
midpoint by $s\in[0,1]$. The midpoint stays fixed, the target gap becomes
$s\Delta$, and the information ratio becomes exactly $s^2K$.

For every induced local policy of a cap-$N$ population algorithm,
\eqref{eq:dual-hellinger-comparison} gives
\[
 h(P_-^A,P_+^A)\le\frac{s^2K}{4}\E_{\bar q}T_A.
\]
This inequality also covers policies that sometimes take zero
observations. Their zero-work event depends only on the seed and has zero
signed contrast; conditioning on activation gives a root-active policy,
and both information and expected work scale by the activation
probability. Suppressing the seed can only decrease information.

Apply Lemma~\ref{lem:tf-midpoint}.
It realizes each hybrid as a one-genuine-unit policy, simulating other
units from independent midpoint tapes in its auxiliary seed. The hybrid
global transcript is a function of that local transcript, so data
processing applies even if the seed reveals the other tapes. Summing the
midpoint local work counts every population observation exactly once.
Therefore
\[
 h_{\mathrm{pop}}\le s^2BK/4.
\]
If $K=0$ take $s=1$; otherwise take $s=\min\{1,(BK)^{-1/2}\}$.
Then $h_{\mathrm{pop}}\le1/4$. Lemma~\ref{lem:affinity-testing} yields
\[
 R_N^{\mathrm{tot}}(B)\ge\frac38s\Delta
 \ge\frac38\frac\Delta{\sqrt{4+BK}}.
\]
Taking the pair supremum proves the first claim; the second follows from
Theorem~\ref{thm:midpoint-modulus}.
\end{proof}

\begin{corollary}[An improved finite characterization]
\label{cor:dual-finite-characterization}
For every integer $N\ge1$,
\begin{equation}
 \left(\frac{P_N(B)}8-
       \frac\rho2\sqrt{\frac B{N+1}}\right)_+
 \le R_{\mu,\rho}(B)\le P_N(B).
 \label{eq:dual-finite-certificate}
\end{equation}
At the explicit cap
\[
 N_*=\left\lceil64\max\{B,\rho B^2\}\right\rceil,
\]
this gives
\begin{equation}
 \boxed{R_{\mu,\rho}(B)\le P_{\mu,\rho,N_*}(B)
                         \le16R_{\mu,\rho}(B).}
 \label{eq:dual-factor-sixteen}
\end{equation}
\end{corollary}

\begin{proof}
Clip any estimator to
$[0,\rho]$ and, before requesting observation $N+1$, returns $\rho/2$.
Markov's inequality and the $L^2$ triangle inequality give
\[
 R_N^{\mathrm{tot}}(B)
 \le R_{\mu,\rho}(B)+\frac\rho2\sqrt{\frac B{N+1}}.
\]
Combining this with Theorem~\ref{thm:dual-adaptive} and the constructive
flow upper bound proves \eqref{eq:dual-finite-certificate}.
Lemma~\ref{lem:tf-floor} gives
\[
 R_{\mu,\rho}(B)\ge\frac1{16}
                 \min\{\rho,\sqrt{\rho/B}\}.
\]
At $N=N_*$ the truncation remainder is at most that floor. Thus
$P_{N_*}(B)\le8(R+R)=16R$.
\end{proof}

The program remains a uniform risk certificate rather than an exact formula
for minimax risk, and its finite dimension does not imply an efficient
algorithm at every budget. Theorem~\ref{thm:entropy-tail-law} supplies a
complementary additive expression involving only the magnitude law, with
its matching lower proof in Section~\ref{sec:matching-lower}. Its
conservative scalar constants are distinct from this program factor.


\begin{remark}[Scope of the factor-sixteen guarantee]
The comparison uses the sufficient cap in
Corollary~\ref{cor:dual-finite-characterization}, or a cap that passes
the verified genie stopping condition in Section~\ref{sec:genie}.
It is not a comparison for every fixed cap. In particular, an ordinary
fixed-cap dual alone is not a lower certificate for unrestricted risk.
With verified bounds $L\le G_N$ and $P_N\le U$, the orientation is
\[
 L\le G_N\le P_\infty\le P_N\le U,\qquad L/8\le R\le U.
\]
Thus $U\le2L$ implies $U\le16R$.
\end{remark}

\section{Finite certificates for the unrestricted experiment}
\label{sec:genie}

The universal cap in Corollary~\ref{cor:dual-finite-characterization} is a
sufficient theoretical cap. This section gives a different computational
route: upper and lower programs at a chosen cap, with an a posteriori
stopping criterion. The lower program strengthens the one-unit experiment
by revealing its latent magnitude and sign if it reaches that cap.
This reveal-at-cap device has a direct precedent in Lee--Valiant
\cite[Section 5.2]{lv}; the contribution here is its convergent robust
dual bracket and exact stopping certificate. The extra information is
used only in a lower-bound argument. The upper
estimator continues to observe ordinary Bernoulli samples alone.

\subsection{The infinite-cap dual and its compactness argument}

Define
\[
 P_\infty(B)=\inf_{N\ge1}P_N(B),\qquad
 K_\infty(q_L,q_H;q_v,q_w)=\sup_{N\ge1}K_N(q_L,q_H;q_v,q_w).
\]
The latter definition uses only deterministically capped local policies.
No assertion about a general unbounded policy is needed for the following
interchange.

\begin{lemma}[The exact dual persists in the limit]\label{lem:infinite-dual}
For every law and every $B>0$,
\begin{equation}\label{eq:infinite-dual}
 P_\infty(B)^2
 =\max_{q_L,q_H,q_v,q_w\in\mathcal Q_\rho}
 \frac{[\theta(q_H)-\theta(q_L)]_+^2}{4+B K_\infty(q_L,q_H;q_v,q_w)}.
\end{equation}
The maximum is attained and is positive. If $D_\infty$ is the corresponding
midpoint-reference modulus, then
\begin{equation}\label{eq:infinite-comparison}
 D_\infty\le P_\infty\le3D_\infty,
 \qquad R(B)\le P_\infty(B)\le8R(B).
\end{equation}
\end{lemma}

\begin{proof}
The space $\mathcal Q_\rho^4$ is weak-star compact. For each fixed cap,
$K_N$ is lower semicontinuous: every fixed policy has a lower semicontinuous
closed-perspective numerator and a positive continuous work denominator,
and taking their supremum preserves lower semicontinuity. Its dual ratio
$f_N=(\Delta\theta)_+^2/(4+BK_N)$ is upper semicontinuous, with value zero
at an infinite denominator. Indeed, the numerator is continuous and
bounded, and lower semicontinuity forces a diverging denominator when
approaching a point with $K_N=\infty$. The functions $f_N$ decrease with $N$.
If $M=\inf_N\max f_N$, the nonempty compact sets $\{f_N\ge M\}$ are
nested. Their intersection gives a quadruple with $\inf_Nf_N\ge M$.
The reverse inequality is immediate. This proves \eqref{eq:infinite-dual}
and attainment. Positivity follows from \eqref{eq:dual-positive-value}.

The equal-third reference absorption in
Theorem~\ref{thm:midpoint-modulus} holds simultaneously at every cap,
hence also for $K_\infty$. It proves the first two comparisons.
For a fixed midpoint pair, the hybrid Hellinger bound holds at every
finite population cap, with the same coefficient $K_\infty$. Increasing
capped prefixes extends the affinity bound to the full terminal
experiment of any finite-expected-work algorithm. The testing and
interpolation argument of Theorem~\ref{thm:dual-adaptive} therefore gives
$R\ge3D_\infty/8$. The constructive inequality $R\le P_N$ holds for
every $N$. These facts prove \eqref{eq:infinite-comparison}.
\end{proof}

\subsection{Revealing the latent state at the cap}

Fix four priors and write $\delta=q_H-q_L$. Continue to use the prefix
coefficients from Proposition~\ref{prop:dual-bellman}. For a prefix
$s=(n,k)$, let
\[
 \ell_s^\pm(d)=(1/2\pm d)^k(1/2\mp d)^{n-k}.
\]
The ordinary stopping reward is
$r_s=\mathfrak r(p_H(s)-p_L(s),p_v(s))$. Define the latent-reveal reward
\begin{equation}\label{eq:genie-terminal}
 g_s=\int
 \left\{\frac{\delta(d)^2\ell_s^+(d)}{q_v(d)}+
             \frac{\delta(d)^2\ell_s^-(d)}{1-q_v(d)}\right\}\,\mu(dd).
\end{equation}
Terms with zero numerator and denominator are zero; a nonzero numerator
over zero has value infinity. In normalized count coordinates, multiply
$r_s,g_s,p_w(s)$ by $\binom nk$ and use the corresponding Pascal factors.

Let $K_N^+$ be the optimal information/work ratio in the experiment which
uses ordinary observations until it stops before depth $N$, but receives
$(D,S)$ at no extra cost if it reaches depth $N$. Its Bellman recursion is
exactly \eqref{eq:dual-bellman}, with terminal value $g_{Nk}$ in place of
$r_{Nk}$. The root condition is unchanged. Thus $\kappa\ge K_N^+$ exactly
when this modified recursion has root surplus at most zero.

For a finite latent divergence, set
\begin{align}
 C_v&=\int\frac{\delta(d)^2}{q_v(d)(1-q_v(d))}\,\mu(dd),\notag\\
 e_N^v&=C_v-
 \sum_{k=0}^N\binom Nk
 \mathfrak r\bigl(p_H(N,k)-p_L(N,k),p_v(N,k)\bigr).
 \label{eq:genie-residual}
\end{align}

\begin{theorem}[Monotone stopping-ratio brackets]\label{thm:genie-ratio}
For every admissible quadruple of priors, with information ratios
interpreted in the extended nonnegative reals,
\begin{equation}\label{eq:genie-bracket}
 K_N\uparrow K_\infty\le K_N^+,
 \qquad K_{N+1}^+\le K_N^+.
\end{equation}
If, in addition, $C_v<\infty$, then
\begin{equation}\label{eq:genie-error}
 0\le K_N^+-K_N\le e_N^v,\qquad e_N^v\longrightarrow0.
\end{equation}
Consequently $K_N^+\downarrow K_\infty$ under this finite-divergence
hypothesis. Without it, the monotone brackets remain valid, but
pointwise convergence of the upper bracket need not hold.
\end{theorem}

\begin{proof}
Under the variance-reference latent law, define the score
\[
 Z(D,S)=
 \begin{cases}
  \delta(D)/q_v(D),&S=1,\\
  -\delta(D)/(1-q_v(D)),&S=0.
 \end{cases}
\]
If $C_v<\infty$, it is square integrable, has mean zero, and has second
moment $C_v$. The score of any observed transcript is its conditional
expectation. Thus its signed-likelihood information is
$\E_v[\E_v(Z\mid\text{transcript})^2]$.

Take any capped ordinary policy with cap larger than $N$ and stop its paid
observations at $T\wedge N$. On reaching $N$, the genie reveals $(D,S)$;
these values allow simulation of every future Bernoulli observation and
therefore of the original terminal transcript. Data processing cannot
decrease information when passing to the stronger genie transcript, while
the work under $q_w$ can only decrease. This proves $K_\infty\le K_N^+$
by taking the supremum over finite ordinary caps. The same simulation
shows that an earlier reveal is stronger than a later reveal, so $K_N^+$
is nonincreasing. The other monotonicity and $K_N\le K_\infty$ are
immediate. These arguments also apply to infinite divergences using
the extended-value definition.

For a cap-$N$ genie policy, erase the reveal. The resulting ordinary policy
has exactly the same paid work. Its lost information is
\[
 \E_v\left[\mathbf1_{\{T=N\}}
                 \Var_v(Z\mid\mathcal F_N)\right]
 \le\E_v\Var_v(Z\mid\mathcal F_N)=e_N^v.
\]
Here $\mathcal F_N$ is the first $N$ observations of the completed
counterfactual stream. An independent policy seed can be included; it
adds no latent information. Since root-active work under $q_w$ is at
least one, dividing by work and taking suprema gives
$K_N^+\le K_N+e_N^v$.

The infinite Bernoulli stream identifies its mean almost surely and hence
identifies $(D,S)$ because $D>0$. The martingale
$\E_v(Z\mid\mathcal F_N)$ therefore converges to $Z$ in $L^2$.
This proves $e_N^v\to0$. Its explicit formula is precisely
\eqref{eq:genie-residual}.
\end{proof}

\begin{remark}[Pointwise convergence can fail at infinite latent divergence]
Let $\mu(dd)=2\,dd$ on $(0,1/2]$, $\rho=1/4$, and take
$q_L=0$, $q_H=q_w=1/4$, and $q_v(d)=1/2-d$. All four priors are
admissible. Here
\[
 C_v=\frac18\int_0^{1/2}\frac{dd}{1/4-d^2}=\infty.
\]
The policy which always reaches depth $N$ and accepts the latent reveal
has infinite information and work $N$, so $K_N^+=\infty$ for every
finite $N$.

Nevertheless $K_\infty\le3/4$. To verify this for every capped policy,
parameterize the latent state by its Bernoulli mean $p$. The variance
reference has density $2(1-p)$ on $(0,1)$, the work reference has density
$3/2$ below $1/2$ and $1/2$ above $1/2$, and the signed endpoint
difference has density $-1/2$ below $1/2$ and $1/2$ above $1/2$.
For an ordered prefix $s$ of length $n\ge1$ with $k$ heads, put
$u_s=\int_0^1p^k(1-p)^{n-k}\,dp$. Then
\[
 |p_H(s)-p_L(s)|\le u_s/2,\qquad
 p_v(s)=\frac{2(n-k+1)}{n+2}u_s,
\]
and hence
\[
 r_s\le\frac{n+2}{8(n-k+1)}u_s\le\frac{3n}{8}u_s.
\]
Sum over stopped prefixes and the independent seed. If $U$ denotes the
auxiliary uniform law of $p$, the information of any capped policy is
at most $(3/8)\E_U T$. The work-reference density dominates half the
uniform density, so this is at most $(3/4)\E_{q_w}T$. Taking suprema
gives the claimed bound. The auxiliary law $U$ need not be admissible;
it is used only for this domination calculation.

This example does not contradict the optimized convergence in the next
theorem, whose proof first regularizes the variance reference.
\end{remark}

\subsection{A convergent lower program and a stopping rule}

Define the strengthened-experiment dual value
\begin{equation}\label{eq:genie-program}
 G_N(B)^2=
 \sup_{q_L,q_H,q_v,q_w\in\mathcal Q_\rho}
 \frac{[\theta(q_H)-\theta(q_L)]_+^2}{4+B K_N^+(q_L,q_H;q_v,q_w)}.
\end{equation}
As before, a ratio with infinite denominator is zero.

\begin{theorem}[A posteriori minimax brackets]\label{thm:genie-program}
For every law, rarity and budget in the model,
\begin{equation}\label{eq:genie-program-convergence}
 0<G_N(B)\uparrow P_\infty(B),\qquad P_N(B)\downarrow P_\infty(B),
\end{equation}
and at every integer $N\ge1$,
\begin{equation}\label{eq:genie-risk-bracket}
 \boxed{\qquad G_N(B)/8\le R_{\mu,\rho}(B)\le P_N(B).\qquad}
\end{equation}
For finite-support laws both programs have finite conic formulations.
The condition $P_N(B)\le2G_N(B)$ eventually holds and certifies
$R(B)\le P_N(B)\le16R(B)$.
\end{theorem}

\begin{proof}
The preceding theorem gives $G_N\le P_\infty$ and monotonicity. The
constant priors $q_L=0,q_H=q_v=\rho$ have finite latent divergence and
give a strictly positive lower bound for every $G_N$.

To prove convergence, fix a quadruple with finite $K_\infty$ and positive
target gap. Put $m=(q_L+q_H)/2$ and regularize only the variance reference:
\[
 q_v^\epsilon=(1-\epsilon)q_v+\epsilon m,
 \qquad 0<\epsilon<1.
\]
The new transcript reference dominates $(1-\epsilon)$ times the old one,
so
\[
 K_\infty(q_L,q_H;q_v^\epsilon,q_w)
 \le\frac{K_\infty(q_L,q_H;q_v,q_w)}{1-\epsilon}.
\]
Also $|\delta|\le2m$ and $|\delta|\le2(1-m)$, so the new latent score
has absolute value at most $2/\epsilon$. Its latent divergence is finite.
Theorem~\ref{thm:genie-ratio} now makes its genie ratio converge to its
ordinary infinite-cap ratio. Keep $\epsilon>0$ fixed throughout
this limit in $N$; no convergence of the unregularized
quadruple is used. Hence
\[
 \lim_NG_N(B)^2\ge
 \frac{(\Delta\theta)^2}{4+B K_\infty/(1-\epsilon)}.
\]
Let $\epsilon\downarrow0$ and maximize using
Lemma~\ref{lem:infinite-dual}. The limits are taken in this order;
no convergence uniform in $\epsilon$ is asserted. This proves
$G_N\uparrow P_\infty$.
Equation~\eqref{eq:infinite-comparison} gives the risk bracket.
Since $P_\infty>0$, the ratio $P_N/G_N$ tends to one. The stated stopping
criterion therefore succeeds at a finite cap, with no claimed uniform
rate for that cap.
\end{proof}

For an explicitly given law $\mu=\sum_iw_i\delta_{d_i}$, the genie dual
is a finite SOCP. In the price formulation, use endpoint variables
$u_H=\alpha q_H$, $u_L=\alpha q_L$, work-reference variables
$e=\eta q_w$, and a variance reference $q_v$. Their rarity and box
constraints are linear perspective constraints. For each atom introduce
\[
 (u_{H,i}-u_{L,i})^2\le q_{v,i}r_i^+,
 \qquad
 (u_{H,i}-u_{L,i})^2\le(1-q_{v,i})r_i^-.
\]
At terminal states replace the ordinary quadratic-over-linear stopping
cone by the linear bound
\[
 X_s\ge\frac14\sum_iw_i
           \{\ell_s^+(d_i)r_i^++\ell_s^-(d_i)r_i^-\}.
\]
Keep the ordinary stopping cones at earlier states and the same linear
continuation/root inequalities. Maximize
$\sum_iw_i(u_{H,i}-u_{L,i})-\alpha^2-\eta B$.
This is exactly the squared value in \eqref{eq:genie-program} after
eliminating the price and scale. It uses $O(N^2+m)$ variables, apart from
the coefficient matrices. For a general measure the result remains a
variational statement; a finite-support computational representation is
not silently assumed for an arbitrary law.

\subsection{Certificates that can be checked independently}

In practice a feasible primal policy gives an upper bound $U\ge R$, and
a feasible genie dual gives a square-root lower bound $L\le G_N$.
Consequently the \emph{verified} inequality $U\le2L$ is sufficient for
a factor-16 risk certificate. An approximate solver status or an
uncertified objective is not sufficient. Section~\ref{sec:computation}
uses exact rational checks for the finite rational examples.

A selected dual quadruple can also yield a stronger direct risk lower
bound than $L/8$. Absorb its two references into its endpoints with equal
thirds, and let $t$ be the new target separation. If a Bellman certificate
gives $\kappa\ge K_N^+$ for the original quadruple, the absorbed midpoint
pair has $K_\infty\le\kappa$. With $x=B\kappa$, affinity testing gives
\begin{equation}\label{eq:direct-genie-risk}
 R(B)\ge\max_{0\le s\le1}
          \frac{st}{2}\left(1-\frac{s^2x}{4}\right)_+.
\end{equation}
The square of this bound is the rational expression
\begin{equation}\label{eq:direct-genie-risk-squared}
 \begin{cases}
 \displaystyle\frac{t^2}{4}(1-x/4)^2,&0\le x\le4/3,\\[4pt]
 \displaystyle\frac{4t^2}{27x},&x\ge4/3.
 \end{cases}
\end{equation}
This supplies an exact squared lower certificate when the inputs and
the certified price are rational.

\section{What a fixed observation cap can identify}
\label{sec:identifiability}

The exact dual also identifies an obstruction that is invisible in a
small-cap numerical upper bound: increasing the number of units cannot
resolve two priors with identical cap-$N$ one-unit distributions.

Let $L=\lceil N/2\rceil$ and define the moment-matching diameter
\begin{equation}\label{eq:cap-diameter}
 \Delta_N=\max\left\{
 \theta(q_H)-\theta(q_L):q_L,q_H\in\mathcal Q_\rho,\
 \int(q_H-q_L)d^{2j+1}\,d\mu=0\ (0\le j<L)
 \right\}.
\end{equation}
Compactness ensures that this maximum is attained.

\begin{theorem}[Exact fixed-cap bias floor]\label{thm:cap-identifiability}
For every magnitude law,
\begin{equation}\label{eq:cap-bias-floor}
 \lim_{B\to\infty}P_N(B)=\Delta_N/2.
\end{equation}
If $\mu$ has exactly $m$ distinct positive support points, then
$\Delta_N=0$ if and only if $N\ge2m-1$.
If $\mu$ has infinite support, then $\Delta_N>0$ for every finite $N$.
\end{theorem}

\begin{proof}
The distribution of $N$ Bernoulli observations from one unit determines
the first $N$ moments of its latent Bernoulli mean. The difference of the
$k$th such moments under two priors is
\[
 \int(q_H-q_L)
       \{(1/2+d)^k-(1/2-d)^k\}\,d\mu.
\]
These polynomials span exactly the odd monomials
$d,d^3,\ldots,d^{2L-1}$. Thus the constraints in
\eqref{eq:cap-diameter} are equivalent to equality of the full depth-$N$
count distributions. Equality at depth $N$ implies equality for every
stopped experiment with that cap, by marginalization and independent
randomization. Conversely, the policy that always observes exactly $N$
samples witnesses a positive information numerator whenever those count
distributions differ. Therefore $K_N=0$ exactly for pairs satisfying
\eqref{eq:cap-diameter}; otherwise $K_N>0$ or $K_N=\infty$.

Apply the decreasing compact-set argument from
Lemma~\ref{lem:infinite-dual} to the dual ratios as $B\to\infty$.
Their pointwise limit is $(\Delta\theta)_+^2/4$ when $K_N=0$, and zero
otherwise. Theorem~\ref{thm:exact-dual} gives
\eqref{eq:cap-bias-floor}.

For the support assertion, consider the finite-dimensional subspace
\[
 \mathcal V_L=\operatorname{span}\{d,d^3,\ldots,d^{2L-1}\}
 \subset L^2(\mu).
\]
If the support has $m\le L$ points, polynomial interpolation gives
$p(d)=dQ(d^2)=1$ on that support, with $\deg Q\le m-1$.
Then $1\in\mathcal V_L$, so moment matching forces zero target gap.

If the support has more than $L$ points, $1\notin\mathcal V_L$.
Indeed, a polynomial $Q$ of degree at most $L-1$ cannot agree with
$x^{-1/2}$ at $L+1$ distinct positive values: otherwise repeated Rolle's
theorem would force the nonvanishing $L$th derivative of $x^{-1/2}$ to
vanish somewhere. The same argument covers infinite support, since an
almost-everywhere polynomial identity extends to the support.

Let $h=1-\operatorname{proj}_{\mathcal V_L}1$. It is a nonzero bounded
function, orthogonal to every required odd monomial, and
\[
 \int h\,d\mu=\int h^2\,d\mu>0.
\]
Set $\epsilon=\rho/(4\|h\|_\infty)$ and
$q_H=\rho/2+\epsilon h$, $q_L=\rho/2-\epsilon h$.
Both priors lie pointwise between $\rho/4$ and $3\rho/4$ and hence
belong to $\mathcal Q_\rho$. They match the odd moments but have positive
target gap $2\epsilon\int h^2\,d\mu$. This proves strict positivity.
Finally $L\ge m$ is equivalent to $N\ge2m-1$.
\end{proof}

The proof also yields the explicit lower estimate
\[
 \Delta_N/2\ge
 \frac{\rho\|h\|_{L^2(\mu)}^2}{4\|h\|_{L^\infty(\mu)}}
 \quad\text{when }1\notin\mathcal V_L.
\]
This theorem identifies the target functional from cap-limited laws;
it does not require recovery of the entire allocation $q$.
This theorem concerns the fixed-cap program and its infinite-budget limit.
It does not impose a bias floor on the unrestricted experiment. The
monotone brackets in Section~\ref{sec:genie} increase the cap when needed,
and their stopping criterion tests the remaining gap directly.

\section{Exact numerical certificates on finite laws}
\label{sec:computation}

The computations in this section produce two-sided certificates for
eight finite rational instances. Numerical conic optimization selects
candidate policies and dual priors. Separate rational-arithmetic
verifiers then check the resulting bounds, including the observation
budget and the stopping criterion. An approximate solver objective is
never used as a certified endpoint.

\subsection{What the two endpoints mean}

Write $U_N$ for the risk bound of a verified observable cap-$N$ policy,
and $L_N$ for the value of a verified feasible genie dual. Their roles are
\[
 L_N\le G_N\le P_\infty\le P_N\le U_N,
 \qquad L_N/8\le R_{\mu,\rho}(B)\le U_N.
\]
Thus $L_N$ is a lower bound on a program value. It is not, by itself,
the reported lower bound on fully adaptive minimax risk. We also compute
the stronger pair-specific adaptive-risk lower bound
$\underline r_N$ from
\eqref{eq:direct-genie-risk}--\eqref{eq:direct-genie-risk-squared}.
Its target separation is that of the absorbed midpoint pair.

The cap search stops only after the exact rational comparison
\begin{equation}\label{eq:computed-stopping-test}
 U_N^2\le4L_N^2
\end{equation}
has passed. This implies $U_N\le16R_{\mu,\rho}(B)$.
The comparison uses the stored rational squared values, not the rounded
decimals in the tables.

\subsection{The benchmark laws and results}

Every instance has $\rho=1/10$. The four laws are
\begin{align*}
 \mu_{\rm one}&=\delta_{1/4},\\
 \mu_{\rm three}&=\tfrac34\delta_{2/5}
                  +\tfrac15\delta_{1/5}+\tfrac1{20}\delta_{2/25},\\
 \mu_{\rm hard}&=\tfrac{99}{100}\delta_{9/20}
                  +\tfrac1{100}\delta_{2/25},\\
 \mu_{\rm label}&=\delta_{1/2}.
\end{align*}
Each was run at $B=100$ and $B=1000$. The candidate caps were
$4,8,16,32$, with the perfect-label cases starting at $N=1$.
Seven instances passed \eqref{eq:computed-stopping-test} at their first
cap. The three-scale law at $B=1000$ failed at $N=4$ and passed at $N=8$.
There were therefore nine attempted brackets and eight successful ones.
The failed bracket is retained in the source bundle; it remains a valid
bracket but does not meet the requested factor-two program criterion.

\begin{table}[htbp]
\centering
\small
\setlength{\tabcolsep}{4pt}
\begin{tabular}{lrr rrr}
\toprule
Law & $B$ & $N$ & $L_N$ & $\underline r_N$ & $U_N$\\
\midrule
One gap & 100 & 4 & 0.03474453 & 0.00619890 & 0.03891866\\
One gap & 1000 & 4 & 0.01461177 & 0.00196026 & 0.01825093\\
Three scales & 100 & 4 & 0.03331182 & 0.00573109 & 0.04551457\\
Three scales & 1000 & 8 & 0.01569090 & 0.00212025 & 0.02553208\\
Rare hard atom & 100 & 4 & 0.02912416 & 0.00459699 & 0.03009854\\
Rare hard atom & 1000 & 4 & 0.01105026 & 0.00145369 & 0.01166586\\
Perfect labels & 100 & 1 & 0.02572478 & 0.00384900 & 0.02572479\\
Perfect labels & 1000 & 1 & 0.00932054 & 0.00121716 & 0.00932055\\
\bottomrule
\end{tabular}
\caption{Successful exact brackets. Both lower columns are rounded
downward and the upper column upward. The adaptive minimax risk lies
in $[\underline r_N,U_N]$; the unrestricted flow-program value lies
in $[L_N,U_N]$. These two intervals have different meanings.}
\label{tab:exact-brackets}
\end{table}

For the failed three-scale attempt at $B=1000,N=4$, the verified
program endpoints, rounded outward, are $0.01359364$ and $0.02989647$,
and the exact ratio exceeds two. At $N=8$ the program ratio is about
$1.6272$, so the exact test passes. Across the eight successful
instances, the program bracket ratios range from essentially one to
$1.6272$. The sharper adaptive-risk bracket ratios range from about
$6.27$ to $12.05$, all below sixteen.

These are small benchmark instances. Their modest successful caps
provide no uniform bound on the cap or runtime needed for another law,
rarity, or budget. The convergence theorem is qualitative and the
certificates concern the particular rational inputs displayed above.

\subsection{An analytic check with perfect labels}

For $\mu=\delta_{1/2}$, the first observation reveals the sign exactly.
Further observations add no information. For every cap $N\ge1$,
\begin{equation}\label{eq:perfect-label-program}
 P_N(B)^2=P_\infty(B)^2=G_N(B)^2
 =\frac{ac}{a+c},
 \qquad
 a=\frac{\rho^2}{4},\quad c=\frac{\rho(1-\rho)}{B}.
\end{equation}
This is an analytic identity for the flow program, without an assertion
that it is the exact fully adaptive minimax risk.

To verify \eqref{eq:perfect-label-program}, conditional averaging removes
randomness in the output given the revealed sign and cannot increase
its second moment. Let $\lambda\le B$ be the active Poisson intensity,
let $y_0,y_1$ be the two outputs, and put
$s=\lambda(y_1-y_0)$. Optimizing the deterministic offset gives worst
bias $\rho|1-s|/2$. At the prior $q=\rho$, centering the outputs minimizes
the raw second moment, giving the variance lower bound
$s^2\rho(1-\rho)/\lambda$. It is attained by
\[
 y_0=-\rho s/\lambda,\qquad y_1=(1-\rho)s/\lambda.
\]
Since $\rho\le1/4$, their second moment is increasing in $q$ on
$[0,\rho]$, so its worst value occurs at $\rho$. Choosing $\lambda=B$
therefore reduces the objective exactly to
$a(1-s)^2+cs^2$, whose minimum is $ac/(a+c)$.
The genie already reveals nothing new after the first observation,
which proves the same equality for $G_N$.

For the two stored perfect-label instances, the exact rational analytic
value lies between the exact rational lower and upper certificates.
This was checked directly, without rounded decimal comparisons.

\subsection{How the certificates are repaired and checked}

The conic solver may return a status indicating limited numerical
accuracy. Such a status neither establishes nor invalidates a certificate.
The following separate checks determine feasibility.

For a lower certificate, the selected priors are quantized to rational
numbers, clipped to $[0,1]$, and, if needed, scaled to satisfy the rarity
constraint exactly. Variance and work references are regularized by
mixing a weight $10^{-8}$ of the constant prior $\rho/2$.
All binomial kernels and latent rewards are then recomputed exactly.
A numerical bisection proposes the work price. A new backward Bellman
pass uses rational arithmetic and rounds each potential upward to a
dyadic grid of denominator $2^{96}$. The price is increased if the
exact root inequality does not hold. An independent verifier reconstructs
the kernels and checks every rarity, stopping, continuation, terminal,
and root inequality.

For an upper certificate, a separate verifier works with the actual
observable policy: its stopping probabilities at each count state,
terminal outputs, Poisson intensity, and deterministic offset. It
computes exact conditional means, second moments, and observation
counts for every atom and sign. Total stopping probability must equal
one exactly. The worst-prior linear expressions are evaluated by an
exact fractional-knapsack calculation over $\mathcal Q_\rho$.
If rationalization raises expected work above $B$, the intensity is
reduced by the exact required factor and outputs are increased by its
reciprocal. This preserves means, and the variance bound is recomputed.
The reported risk upper is
\[
 U_N^2=
 \left(\sup_{q\in\mathcal Q_\rho}
                  |\operatorname{Bias}_q|\right)^2+
 \sup_{q\in\mathcal Q_\rho}\operatorname{Var}_q.
\]
The two maxima need not occur at the same prior; their sum remains a
valid upper bound. Integer square-root calculations round every
displayed lower RMS downward and every displayed upper RMS upward.
The verifier also checks the direction of the stored decimal rounding.

In addition to the optimized genie brackets, the bundle includes eight
ordinary finite-cap dual checks and sixteen fixed-quartet midpoint-genie
checks. Standard-library-only re-verification accepted all 33 stored
lower certificates and nine upper policies. Six targeted corruptions
were rejected: a rarity violation, an insufficient Bellman price, an
overstated lower value, an insufficient genie terminal reward, excessive
actual policy work, and an understated upper risk. These checks remain
active when Python is invoked with its optimization flag.

\subsection{Standalone reproduction}

The source directory \texttt{checks/} is self-contained. It includes
a clearly attributed copy of the unchanged round-three primal solver,
the independent dual solver, the exact verifiers, and all certificate
files. It does not import a solver from a neighboring research directory.
From \texttt{checks/}, verification needs only Python's standard library:
\begin{verbatim}
python self_check_certificates.py
python verify_finite_dual.py \
  optimized_genie/three_scales_N8_B1000_lower.json
python verify_policy_upper.py \
  optimized_genie/three_scales_N8_B1000_upper.json
\end{verbatim}
The numerical search used Python 3.12 with CVXPY 1.9.3, NumPy 2.5.4,
SciPy 1.18.1, and Clarabel 0.11.1. These observed versions are pinned
in \texttt{requirements.txt}. The principal search is reproduced by
\begin{verbatim}
python -m pip install -r requirements.txt
python solve_genie_bracket.py
\end{verbatim}
Re-running a numerical solver may select different feasible priors
and slightly different bounds. The mathematical guarantee comes from
the independently checked rational artifacts.

\clearpage
\begin{thebibliography}{99}
\bibitem{original}
Supplied manuscript. \emph{A Distributional Minimax Law for Adaptive
Threshold Estimation}, version 4, 26 October 2026.
User-supplied \texttt{paper.pdf}.

\bibitem{repo}
{\raggedright
\emph{adaptive-threshold-estimation}, version-5 project baseline.
Pinned commit:\par
\texttt{eabf274cb42d4bff4daab6b58568ad6169e7f0a7}.\par
\href{https://github.com/davidiach/adaptive-threshold-estimation/tree/eabf274cb42d4bff4daab6b58568ad6169e7f0a7/paper/v5}{Pinned version-5 manuscript source}.\par}

\bibitem{lv}
J. C. H. Lee and P. Valiant.
\emph{Uncertainty about Uncertainty: Optimal Adaptive Algorithms for
Estimating Mixtures of Unknown Coins}.
Proceedings of SODA, 414--433, 2021.
Extended version arXiv:1904.09228v3.
\url{https://arxiv.org/html/1904.09228v3}.

\bibitem{cl16}
L. Chen and J. Li.
\emph{Open Problem: Best Arm Identification: Almost Instance-Wise
Optimality and the Gap Entropy Conjecture}.
Proceedings of COLT, PMLR 49:1643--1646, 2016.
\url{https://proceedings.mlr.press/v49/chen16b.html}.

\bibitem{clq}
L. Chen, J. Li, and M. Qiao.
\emph{Towards Instance Optimal Bounds for Best Arm Identification}.
Proceedings of COLT, PMLR 65:535--592, 2017.
\url{https://proceedings.mlr.press/v65/chen17b.html}.

\bibitem{jn}
A. B. Juditsky and A. S. Nemirovski.
\emph{Nonparametric estimation by convex programming}.
The Annals of Statistics 37(5A):2278--2300, 2009.
\url{https://arxiv.org/abs/0908.3108}.

\bibitem{pw}
Y. Polyanskiy and Y. Wu.
\emph{Dualizing Le Cam's method for functional estimation, with applications
to estimating the unseens}. arXiv:1902.05616v2, 2021.
\url{https://arxiv.org/abs/1902.05616v2}.

\bibitem{bv}
S. Boyd and L. Vandenberghe. \emph{Convex Optimization}.
Cambridge University Press, 2004. Chapter 5.
\url{https://web.stanford.edu/~boyd/cvxbook/}.

\bibitem{cs}
M. J. Cho and R. H. Stockbridge.
\emph{Linear Programming Formulation for Optimal Stopping Problems}.
SIAM Journal on Control and Optimization 40(6):1965--1982, 2002.
\url{https://doi.org/10.1137/S0363012900377663}.

\bibitem{hoeffding}
W. Hoeffding. \emph{Probability Inequalities for Sums of Bounded Random
Variables}. Journal of the American Statistical Association 58(301):13--30,
1963. \url{https://doi.org/10.1080/01621459.1963.10500830}.

\bibitem{sion}
M. Sion. \emph{On general minimax theorems}.
Pacific Journal of Mathematics 8(1):171--176, 1958.
\url{https://doi.org/10.2140/pjm.1958.8.171}.

\bibitem{cvxpy}
CVXPY documentation. \emph{Second-order cone program}.
\url{https://www.cvxpy.org/examples/basic/socp.html}.
The searches used CVXPY 1.9.3 and Clarabel 0.11.1;
the exact certificate checkers use the Python standard library.
\end{thebibliography}
\end{document}
